\documentclass[a4paper]{article}
\usepackage{luatexja}
\usepackage{amsmath,amssymb,amsfonts}
\usepackage{amsbsy}
\usepackage{amsthm}
\usepackage{geometry}
\usepackage{enumitem}
\usepackage{booktabs}
\usepackage{hyperref}

\geometry{margin=1in}

\newtheorem{theorem}{定理}[section]
\newtheorem{lemma}[theorem]{補題}
\newtheorem{proposition}[theorem]{命題}
\newtheorem{corollary}[theorem]{系}
\theoremstyle{definition}
\newtheorem{definition}[theorem]{定義}
\newtheorem{example}[theorem]{例}
\theoremstyle{remark}
\newtheorem{remark}[theorem]{注意}
\renewcommand{\proofname}{証明}

\newcommand{\Jnorm}{J_{\rm norm}}
\newcommand{\Jbound}{\frac{3\pi^2}{8}}

\title{離散双四元数二重時空：ねじれ不整合およびPGA埋め込みによるディオファントス予想の厳密証明}

\author{https://github.com/hypernumbernet}
\date{\today}

\begin{document}

\maketitle

\begin{abstract}
我々は、古典重力とディオファントス数論の双方が単一の構造——離散双四元数二重時空代数——から現れる厳密な代数的枠組みを展開する。16次元双四元数代数 \(\mathbb{H} \otimes_{\mathbb{R}} \mathbb{H} \cong \mathrm{Cl}(3,1)\) を射影幾何代数 \(G(3,1,1)\) に埋め込み、双対時空のラピディティを有限トーラス \((\mathbb{Z}/N\mathbb{Z})^6\) へ離散化することにより、許容構成上で正規化ねじれスカラーが厳密に有界 \(|\Jnorm| \leq 1\)（同値に \(|J| \leq 3\pi^2/8\)）である理論を得る。この上限はキリング二次形式に対する初等的評価と非負ラピディティおよび反同期から従う。

整数は標準ローターおよび零トランスレータを通じて代数へ忠実に埋め込まれ、すべてのディオファントス方程式は有界格子上の有限探索へと変換される。第5章はフェルマー三つ組の二軸幾何座席を展開する。混合種は加法的零軸を時間的零生成子へ漏らし、次数 \(n\ge 3\) の仮の正値解の上ではその流出は有限時間まで生き残り、\(N_1\) の調和振動として時間的零成分の非零倍を残す。混合モーターと冪和との両立不可能性は未閉鎖であり、古典的なフェルマーの最終定理は主張しない。第6--10章は、ビール 予想、abc 予想、リーマン 仮説、ゴールドバッハ 予想、ポリニャック 予想、および コラッツ 予想に対する共通の増幅プログラムを概説する。それらのプログラムを無条件の定理へと転じる問題固有の橋渡しは確立されておらず、それら予想の無条件の主張は述べない。

本枠組みは連続極限 \(N \to \infty\) において一般相対性理論の遠平行等価理論（TEGR）を正確に回復する一方、離散構造は数論的基礎と、特異点および発散を除去する自然な紫外カットオフの双方を供給する。したがって同一の代数的対象が、追加の場、次元、あるいは量子化の要請なしに算術と重力を統一する。

本理論は具体的な物理的含意をもたらす。離散的な重力波エコー、各コンパクト化 \(N\) における有限個の許容核ローター類、および核エネルギー準位への離散的なねじれ補正である。数値的カットオフ \(A_* \lesssim 300\) は未導出の発見的主張であり、定理ではない。

この再定式化において、時空は離散的、粒子局所的、かつ二重時間的である。重力は、すべての質量粒子が運ぶ内在的な双対時空の間のねじれ不整合の集団的知覚として再解釈される。
\end{abstract}

\section{はじめに}

二重時空理論は、すべての質量粒子が運ぶ内在的に対をなす2つの ミンコフスキー 構造の間のねじれ不整合として重力と慣性を符号化し、これを16実次元双四元数代数 \(\mathbb{H} \otimes_{\mathbb{R}} \mathbb{H} \cong \mathrm{Cl}(3,1)\) の内部で実現する。パラメータ空間のねじれ不整合スカラー
\[
J = \frac12\sum_{a=1}^3(\alpha_a^2-\beta_a^2),
\]
は、6生成子の ローレンツセクター \(\mathfrak{so}(3,1)\)（12次元の和 \(\mathfrak{so}(3,1)\oplus\mathfrak{so}(3,1)\) ではない）から構成され、ローターからテトラッドへの辞書が固定された後にのみ遠平行スカラーと同一視される。アインシュタイン--ヒルベルト作用との変分的同値は TEGR 文献から取り、ここでは再導出しない。半角展開のもとでは \(B(\Omega_{\rm biv},\Omega_{\rm biv})=2\sum_a(\alpha_a^2-\beta_a^2)\) であるから、\(J\) は \(\frac1{16}B(\Omega_{\rm biv},\Omega_{\rm biv})\) の4倍である。

枠組みを厳密な数論的設定へ引き上げるため、二つの決定的な拡張を導入する。第一に、双対写像 \(X \mapsto Xi\) により誘導される自然なコンパクト性を、ラピディティパラメータの離散化によって形式化する：
\[
\omega_a = \frac{2\pi n_a}{N}, \qquad \phi_a = \frac{2\pi m_a}{N}, \quad n_a,m_a \in \mathbb{Z},
\]
ここで \(N \gg 1\) は有効コンパクト化パラメータである。得られる位相空間は有限トーラス \((\mathbb{Z}/N\mathbb{Z})^6\) である。有限なのはこのトーラスの離散ローター像であり、未指定の整数双四元数次数の単数群ではない。許容錐に制限すれば、ディオファントス問題はその像上の有限探索となる。

第二に、零ベクトル \(e_4\)（\(e_4^2 = 0\)）を添加することにより、双四元数代数を射影幾何代数 \(G(3,1,1)\) へ埋め込む。これにより、3個の双曲型生成子、3個の巡回生成子、および4個の零生成子 \(N_\mu = e_4 \wedge e_\mu\) からなる10次元の二重ベクトル空間が得られる。零セクターは並進的（加法的）自由度の標準的な代数的実現を供給し、拡張不変量 \(J^{(5)}\) は非退化のままである。さらに \(G(3,1,1)\) の平面ベースの括弧積は、双対時空を3-ブレードへの法バンドルと同一視し、加法構造と乗法構造の間の幾何学的橋渡しを与える。

通常セクター・ローター上のダガー演算 \(R_{\rm usual}^\dagger\) は、結合不整合から双対セクターのデータを抽出し（同値に、ラピディティ対 \((\alpha_a,\beta_a)\mapsto(\beta_a,\alpha_a)\) を入れ替える）、続く双対ローターによる再構成は、\(J\) の符号を除いてねじれ構成を回復する。離散ディオファントス問題の決定可能性は、許容トーラスの有限性と、網羅的格子探索または \(|\Jnorm|\le 1\) に対する増幅による矛盾（第10章）のいずれかから従う。ダガーそれ自体はねじれ高さを減少させない。

この統一された代数的枠組みの内部で、七つの古典問題は二つの異なる座席を占める。フェルマーの最終定理は、射影三つ組 \([a:b:c]\) の二軸ねじれ種に結び付く（第5章）。次数 \(n=2\) は純巡回、次数 \(n\ge 3\) は混合であり、混合種はスカラー \(J\) には見えない異軸干渉を運ぶ。すでに一階で、それらは加法的零軸 \(N_1\) を時間的零生成子 \(N_0\) へ漏らす。仮の正値解の上では特性は楕円的であり、流出は有限時間まで \(N_1\) の調和サンドイッチとして生き残り、その剰余は \(\alpha>0\) のときに限り \(N_0\) の非零倍である。生きている残差は、その剰余を正値冪和との両立不可能性へ昇格させることである。それが閉じるまで、古典的なフェルマーの最終定理は主張しない。ビール 予想、abc 予想、リーマン 仮説、ゴールドバッハ 予想、ポリニャック 予想、および コラッツ 予想は、有限トーラス上の増幅プログラムのままである。それぞれ不等指数の不整合、品質上限、アンサンブル、分解、偶数間隙リセット、または \(J\) を最小化しようとする流れへと還元される。それらの橋渡しは未証明のままであり、六つの予想の無条件の主張は述べない。二重時空四値論理（D4L）は、その禁止をレジーム計算として記録する。証明済みの核、棄却済み診断、未閉鎖の残差は二値フラグメントが同時にホストできず、どれも古典予想を \(T\)-含意しない（\(\texttt{dst-4-valued-logic-ja.tex}\)）。

以降の許容解析は、許容構成上の厳密な有界性 \(|\Jnorm| \leq 1\)（定理~\ref{thm:torsion-bound}）に依拠する。これはキリング形式の符号非対称性と、非負の埋め込みラピディティおよび反同期 \(\alpha_a + \beta_a \leq \pi/2\) から導かれる。

本論文の構成は以下のとおりである。第2章は離散二重時空と \(G(3,1,1)\) 埋め込みを導入する。第3章はねじれ有界性の厳密定理を確立する。第4章は数論の離散双四元数環への代数的埋め込みを展開する。第5章はフェルマーの二軸プログラムを展開し、混合性が整数軸の有限時間の性質であることを示す（混合モーター残差は未閉鎖である）。第6章はビール プログラム、第7章は abc プログラム、第8章はリーマン プログラム、第9章は残る三つの予想を記録する。第10章は統一ディオファントス降下枠組みを提示する。第11章は物理的および数論的含意を論じ、第12章は今後の方向で結ぶ。

本研究は、離散性と射影幾何構造を備えた双四元数代数が、古典重力と算術数論の双方に対する完全な代数的基礎を供給することを示す。

\section{離散二重時空と PGA \texorpdfstring{\(G(3,1,1)\)}{G(3,1,1)} 埋め込み}

\subsection{双対時空の離散化}

双四元数代数の双対写像 \( X \mapsto Xi \) は、ラピディティパラメータ上に内在的周期性を誘導する。したがって完全ローターの六つの実パラメータを
\[
\omega_a = \frac{2\pi n_a}{N}, \qquad \phi_a = \frac{2\pi m_a}{N}, \quad a=1,2,3,
\]
に従って離散化する。ここで \( n_a, m_a \in \mathbb{Z} \) であり、\( N \gg 1 \) はプランク単位での粒子質量によって定まるコンパクト化パラメータである。相対角の位相空間はこれにより有限トーラス \( (\mathbb{Z}/N\mathbb{Z})^6 \) へ還元される。その結果、ねじれ不整合ローターは \(\mathrm{Spin}^+(3,1)\) 内部の有限な離散ローター像に値をとる。その像が有限であるのは定義域が有限だからであり、無限でありうる整数双四元数次数の単数群ではない。この有限性は、連続 ローレンツ構造を極限 \( N \to \infty \) で保ちつつ、すべてのディオファントス問題を離散格子上の有限探索へ変換する。許容性はその格子上の追加の仮説のままである。

\subsection{射影幾何代数 \texorpdfstring{\(G(3,1,1)\)}{G(3,1,1)} への埋め込み}

並進自由度を代数的に取り込むため、双四元数構造を五次元射影幾何代数 \( G(3,1,1) \) へ持ち上げる。\(\{e_0 = j, e_1 = kI, e_2 = kJ, e_3 = kK\}\) を \(\mathrm{Cl}(3,1)\) の標準基底とする。単一の零ベクトル \( e_4 \) を添加し、
\[
e_4^2 = 0, \qquad \{e_4, e_\mu\} = 0 \quad (\mu = 0,1,2,3)
\]
を満たすものとする。得られる代数は32次元である。その10次元二重ベクトル空間は三つの異なる類に分解する：
\begin{itemize}
  \item 双曲型生成子（平方が \( +1 \)）：\( iI, iJ, iK \)、
  \item 巡回生成子（平方が \( -1 \)）：\( I, J, K \)、
  \item 零生成子：\( N_\mu = e_4 \wedge e_\mu \)（\( \mu = 0,1,2,3 \)）。
\end{itemize}
四つの零生成子は強い消滅性 \( N_\mu N_\nu = 0 \) をすべての \(\mu,\nu\) について満たし、各々は反転奇である：\( \widetilde{N_\mu} = -N_\mu \)。交換子 \([X,Y]=XY-YX\) のもとで10次元スパンは \( G(3,1,1) \) の リー部分代数である：六つの双曲型および巡回生成子は自らのスパン上で閉じ、ねじれ生成子と零生成子の括弧はつねに再び零であり、零生成子同士は可換である。したがって零スパンは ローレンツセクターがそのベクトル表現によって作用するアーベルイデアルであり、10次元空間は ポアンカレ 代数 \(\mathfrak{iso}(3,1)\) の半直積構造 \(\mathfrak{so}(3,1)\ltimes\mathbb{R}^{3,1}\) を、次元の数え上げではなく代数的事実として担う。ローレンツセクター内部の構造定数は標準形 \([B^-_a,B^-_{a+1}]=2B^-_{a+2}\)、\([B^+_a,B^-_{a+1}]=2B^+_{a+2}\)、\([B^+_a,B^+_{a+1}]=-2B^-_{a+2}\) をとり、最後の符号こそが \(\mathfrak{so}(3,1)\) をコンパクトな \(\mathfrak{so}(4)\) から分かつものである。

\subsection{拡張ねじれ不整合と不変量 \texorpdfstring{\( J^{(5)} \)}{J}}

拡張モーターの無限小生成子は
\[
\Omega_{\rm biv}^{(5)} = \underbrace{\sum_{a=1}^3 \Bigl( \frac{\alpha_a}{2} B_a^+ + \frac{\beta_a}{2} B_a^- \Bigr)}_{\Omega_{\rm torsion}} + \underbrace{\sum_{\mu=0}^3 \frac{\lambda^\mu}{2} N_\mu}_{\Omega_{\rm trans}}
\]
と分解される。ここで \( B_a^+ \) および \( B_a^- \) はそれぞれ双曲型および巡回生成子を表す。すべての零二重ベクトルは平方が零であり自身と可換であるため、並進部分の指数は一次で打ち切られる：
\[
\exp(\Omega_{\rm trans}) = 1 + \sum_{\mu=0}^3 \frac{\lambda^\mu}{2} N_\mu.
\]
モーターは順序積 \( M := RT \) によって定義され、\( R=\exp(\Omega_{\rm torsion}) \) および \( T=1+\Omega_{\rm trans} \) である。両因子はユニタリであり、したがって \( M\widetilde{M}=1 \) が成り立つ。順序は無害ではない：\([\Omega_{\rm torsion},\Omega_{\rm trans}]\neq 0\) のとき、積 \( RT \) は \( TR \) とも \( \exp(\Omega_{\rm biv}^{(5)}) \) とも異なる。この選択を無害にするのは、上で記した零セクターのイデアル構造を、リー代数から群へ持ち上げたものである。指数による共役は随伴作用の指数であり、\( e^{\Omega}\,X\,e^{-\Omega} = e^{\operatorname{ad}\Omega}X \) が成り立つ。したがって \( \operatorname{ad}\Omega \) で安定な部分空間はすべて、サンドイッチ \( e^{\Omega}\,(\cdot)\,\widetilde{e^{\Omega}} \) のもとでも安定である。ローレンツセクターは零スパンを自らの中へ写すから、あらゆるねじれローターはトランスレータを正規化する：
\[
R\,T\,\widetilde{R} = T' = 1+\Omega'_{\rm trans}.
\]
これは共役された係数を閉じた形で書き下せる純ブーストに限らず、すべての \( R=\exp(\Omega_{\rm torsion}) \) について成立する。さらにユニタリなサンドイッチは代数自己同型であるから指数と交換し、\( R\,e^{\Omega}\,\widetilde R = e^{R\Omega\widetilde R} \) が成り立つ。すなわちねじれローターはねじれローターをねじれローターへ共役する。したがってモーターは半直積の形で掛け合わされる：
\[
(R_pT_p)(R_qT_q) = (R_pR_q)\,T', \qquad T' = \bigl(\widetilde{R_q}\,T_p\,R_q\bigr)\,T_q.
\]
ここでトランスレータはアーベル正規部分群をなし、ローターはその上にベクトル表現を通じて作用する。これを反復すれば \( (RT)^p = R^p\,T_p \)（\( T_p \) は再びトランスレータ）であり、整数冪のもとで増幅するのはねじれローターのみで、並進部分は単に蓄積されるにすぎない。こうしてモーターが生成する群は、生成子空間の半直積構造 \(\mathfrak{so}(3,1)\ltimes\mathbb{R}^{3,1}\) を指数化したものにほかならない。この描像が与えないのは、ローター因子 \( R_pR_q \) をねじれパラメータで表す閉じた公式だけであり、それにはローレンツセクターの指数写像がローター群の上への写像であることが要る。

同じ機構が \( RT \) と真の指数との関係を決着させる。零スパンがアーベルイデアルであることから、全二重ベクトルの指数は
\[
\exp(\Omega_{\rm torsion}+\Omega_{\rm trans}) = R\,(1+Z), \qquad Z = \int_0^1 e^{-s\Omega_{\rm torsion}}\,\Omega_{\rm trans}\,e^{s\Omega_{\rm torsion}}\,ds
\]
と因子分解され、着衣された並進 \( Z \) は再び零スパンに属する。ゆえに \( \exp(\Omega_{\rm biv}^{(5)}) \) はそれ自身ユニタリなモーターであり、\emph{同じ}ねじれローター \( R \) をもち、\( RT \) との違いは並進係数を \( \Omega_{\rm trans} \) からではなく \( Z \) から読み取る点にのみある。ねじれと並進が可換なら \( Z=\Omega_{\rm trans} \) となって二つの規約は一致する。一般には一致しない（動径ブーストと時間並進の組み合わせですでに \( Z\neq\Omega_{\rm trans} \) となる）が、経路 \( t\mapsto e^{t\Omega_{\rm torsion}}(1+t\Omega_{\rm trans}) \) に沿った \( e^{t\Omega_{\rm biv}^{(5)}} \) とのずれは二次で初めて現れ、それはちょうど ベイカー–キャンベル–ハウスドルフ 項 \( \tfrac{t^2}{2}[\Omega_{\rm torsion},\Omega_{\rm trans}] \) である。この項を指数に加えれば両経路は三次まで一致する。すなわちこの不一致は並進イデアル上の座標の取り替えにすぎない。ローター因子、ひいては本論文のあらゆるねじれ量、とくに \( J \) は、\( RT \) と \( \exp(\Omega_{\rm biv}^{(5)}) \) のどちらを選ぶかに依存しない。順序積が選ばれるのは、それが並進係数 \( \lambda^\mu \) を直接に露わにするからである。付随する非退化二次不変量は
\[
J^{(5)} = \frac{1}{16} B_{\rm Killing}(\Omega_{\rm torsion}, \Omega_{\rm torsion}) + \frac{1}{2} \eta_{\mu\nu} \lambda^\mu \lambda^\nu
\]
である。半角展開のもとで第一項は \(\frac18\sum_a(\alpha_a^2-\beta_a^2)\) に等しい。本論文のパラメータ空間スカラーは \( J=\frac12\sum_a(\alpha_a^2-\beta_a^2) \) であり、その項の4倍である。第二項は、光円錐上でちょうど消える並進不整合のローレンツ不変測度を供給する。

\subsection{双対共変性}

親双四元数代数の双対写像 \( X \mapsto Xi \) は擬スカラー \( i=e_0e_1e_2e_3 \) による右乗法であり、\( i^2=-1 \) を満たし、零ベクトル \( e_4 \) と可換であり、ねじれセクターのすべての二重ベクトルと可換である。したがって六次元のねじれ（双曲型--巡回）セクター上で双対性は複素構造であり、生成子上では厳密である：\( B^-_a\,i = B^+_a \) かつ \( B^+_a\,i=-B^-_a \)。すべてのブーストは対応する回転の双対であり、ローレンツセクターは回転代数の複素化であり、括弧は複素双線型 \([Xi,Y]=[X,Y]\,i\)、\([Xi,Yi]=-[X,Y]\) である。上に記録したブースト--ブーストの符号は \( i^2=-1 \) にほかならない。パラメータチャート上では双対性は四分の一回転 \((\alpha_a,\beta_a)\mapsto(\beta_a,-\alpha_a)\) であり、\(J\) を反転し、符号なしの質量 \(\tfrac12\sum_a(\alpha_a^2+\beta_a^2)\) を保つ。各零生成子は次数4の元へ送られるため、10次元生成子空間は双対性のもとで\emph{閉じていない}。並進パラメータは別途追跡されなければならない。通常セクター・ローター上のダガー演算 \( R_{\rm usual}^\dagger \) はブーストパラメータを双対セクターの回転パートナーへ写し、パラメータチャート上では単に \((\alpha_a,\beta_a)\mapsto(\beta_a,\alpha_a)\) を入れ替える。これは擬スカラーによる四分の一回転と、新しい回転パラメータの符号においてのみ異なる。両者はともに \(J\mapsto -J\) を送り、したがって正規化高さ \(|\Jnorm|\) を保存する。トーラスの有限な離散ローター像は、許容性とあわせて、第10章で用いられる探索空間を区切る。ディオファントス論証の停止は、その有限性（および定理~\ref{thm:torsion-bound} に対する増幅）に依拠し、ダガーのもとでの高さの厳密な減少には依拠しない。

したがって本章の構成は、乗法的（ローター）構造と加法的（零トランスレータ）構造が共存し、許容構成上で正規化ねじれスカラーが \(|\Jnorm|\leq 1\) を満たし、すべてのディオファントス方程式が有限格子探索へ還元される、完全な代数的環境を供給する。

\section{ねじれ有界性の厳密定理}

我々は\emph{許容}構成上のねじれ不整合スカラーに対する厳密な上限を導く。上限は初等的である（キリング形式の符号非対称性と、非負ラピディティおよび反同期）。それは Lean~4（プロジェクト \texttt{DstDiophantine}）において機械検証済みである。

\begin{definition}[許容ねじれ構成]
\label{def:admissible}
離散トーラス上で、ラピディティは \(\alpha_a = 2\pi n_a/N\) および \(\beta_a = 2\pi m_a/N\) として埋め込まれ、\(n_a,m_a \in \{0,\dots,N-1\}\) であるから、\(\alpha_a,\beta_a \in [0,2\pi)\) である。構成が\emph{許容}であるとは、
\[
\alpha_a \geq 0, \qquad \beta_a \geq 0, \qquad \alpha_a + \beta_a \leq \frac{\pi}{2}
\quad \text{for all } a = 1,2,3.
\]
を満たすことである。
\end{definition}

\begin{definition}[正規化ねじれスカラー]
\label{def:Jnorm}
\emph{正規化}ねじれスカラーは
\[
\Jnorm := \frac{8}{3\pi^2}\,J
= \frac{4}{3\pi^2}\sum_{a=1}^3 (\alpha_a^2 - \beta_a^2)
\]
である。
\end{definition}

\begin{remark}
主枝条件 \(|\delta_a| = |\alpha_a + \beta_a| \leq \pi/2\) だけでは \(J\) を\emph{押さえない}。例えば \(\alpha_a = 10\)、\(\beta_a = -10 + \pi/4\)（一軸のみ、他は零）は \(|\delta_a| \leq \pi/2\) を満たすが \(|\Jnorm| > 1\) を与える。許容性は追加で \(\alpha_a,\beta_a \geq 0\) を要求する。これは離散埋め込みでは自動的に成り立つが、以下の上限には本質的である。
\end{remark}

\subsection{キリング形式からの導出}

ねじれ不整合二重ベクトルを、ローレンツセクター \( \mathfrak{so}(3,1) \) の標準6生成子基底で表す：
\[
\Omega_{\rm biv} = \sum_{a=1}^3 \Bigl( \frac{\alpha_a}{2} i\Gamma_a + \frac{\beta_a}{2} \Gamma_a \Bigr),
\]
ここで \( \Gamma_1 = I \)、\( \Gamma_2 = J \)、\( \Gamma_3 = K \) である。各コピーの \(\mathfrak{so}(3,1)\) 上のキリング形式は
\[
B(i\Gamma_a, i\Gamma_a) = +8, \qquad B(\Gamma_a, \Gamma_a) = -8
\]
と評価される。したがって生成子展開は
\[
B(\Omega_{\rm biv}, \Omega_{\rm biv}) = 2 \sum_{a=1}^3 (\alpha_a^2 - \beta_a^2)
\]
を与え、同一の展開に対してときに書かれる因子 \(8\sum\) ではない。本論文を通じてねじれスカラーはパラメータ空間の二次形式
\[
J = \frac12 \sum_{a=1}^3 (\alpha_a^2 - \beta_a^2)
\]
であり、上記の展開のもとで \(\frac1{16}B(\Omega_{\rm biv},\Omega_{\rm biv})\) の4倍である。以下の有界性の主張はこの \(J\) および \(\Jnorm=\frac8{3\pi^2}J\) を指す。

\subsection{許容構成上の初等的上限}

\begin{theorem}[ねじれ有界性]
\label{thm:torsion-bound}
すべての許容構成（定義~\ref{def:admissible}）に対して、
\[
|J| \leq \Jbound, \qquad |\Jnorm| \leq 1.
\]
さらに \(|\Jnorm| = 1\) であるのは、すべての軸が同一の極値型であるときに限る。すなわち、すべての \(a\) に対する純双曲型支配 \(\alpha_a = \pi/2\)、\(\beta_a = 0\)（このとき \(J = \Jbound\)）、またはすべての \(a\) に対する純楕円型支配 \(\alpha_a = 0\)、\(\beta_a = \pi/2\)（このとき \(J = -\Jbound\)）である。混合極値（一軸が双曲型で他軸が楕円型）は \(|\Jnorm| < 1\) を与える。許容錐上で \(\Jnorm\) の像はちょうど \([-1,1]\) である。
\end{theorem}

\begin{proof}
許容構成上では、\(\alpha_a + \beta_a \leq \pi/2\) および \(\alpha_a,\beta_a \geq 0\) から \(\alpha_a,\beta_a \leq \pi/2\) が従う。\(\alpha_a^2 - \beta_a^2 = (\alpha_a - \beta_a)(\alpha_a + \beta_a)\) を用いると、各軸について
\[
|\alpha_a^2 - \beta_a^2| \leq (\pi/2)^2
\]
を得る。等号が成立するのは \(\{\alpha_a,\beta_a\}=\{\pi/2,0\}\) のときに限る。\(a = 1,2,3\) について和をとり2で割ると \(|J| \leq \frac12 \cdot 3(\pi/2)^2 = \Jbound\) を得る。\(8/(3\pi^2)\) を掛けると \(|\Jnorm| \leq 1\) となる。

\(|\Jnorm|=1\) ならば、すべての軸が軸ごとの上限を飽和し、三つの符号付き寄与 \(\alpha_a^2-\beta_a^2\) は符号を混ぜられない（さもなくば三角不等式は厳密である）。したがって三軸は同時に双曲型であるか、同時に楕円型である。逆にこれら二つの構成は \(\Jnorm=\pm 1\) を達成する。像に関する主張は、純双曲型光線 \(\alpha_a=t\cdot\pi/2\)、\(\beta_a=0\)（\(t\in[0,1]\)）および純楕円型光線をスケールすることにより従う。
\end{proof}

\subsection{双対写像と Spin 被覆からの制約}

双対写像 \( X \mapsto Xi \) は反同期 \(\alpha_a \approx -\beta_a\) を動機づける。許容離散構成上では、偏差 \(\delta_a = \alpha_a + \beta_a\) は単に \(|\delta_a| \leq \pi/2\) であるのではなく \(\delta_a \leq \pi/2\) を満たす。Spin 群の二重被覆および主枝対数は、ローター冪増幅（付録）で用いられる追加の\emph{物理的}制約を供給する。それらは許容性なしにはそれ自体で \(|\Jnorm| \leq 1\) を含意しない。

\subsection{離散トーラス上の有界性}

ラピディティが \((\mathbb{Z}/N\mathbb{Z})^6\) へ離散化されるとき、許容不整合二重ベクトルの集合は有限である。定理~\ref{thm:torsion-bound} はすべての許容格子点に適用され、上限の格子形はより鋭い。\(k=\lfloor N/4\rfloor\) と書くと、
\[
|\Jnorm| \leq \Bigl(\frac{4k}{N}\Bigr)^2 \leq 1.
\]
\(4\mid N\) ならば \(k=N/4\) であり、上限 \(\pm 1\) は \(n_a=N/4\)、\(m_a=0\)（すべての \(a\)）および \(n_a=0\)、\(m_a=N/4\)（すべての \(a\)）の格子点で達成される。\(4\nmid N\) ならば、すべての許容点は厳密不等式 \(|\Jnorm|<1\) を満たす。\(N\to\infty\) とすると、到達可能な許容値は \([-1,1]\) において稠密である。全トーラス上の非許容点は \(|J| > \Jbound\) でありうる。ディオファントス論証は有限な離散ローター像の内部の許容構成に制限する。

拡張不変量 \(J^{(5)}\) は並進項 \(\frac12 \eta_{\mu\nu}\lambda^\mu\lambda^\nu\) を含み、これは \(\lambda^\mu\) に対する追加の制約なしには\emph{有界ではない}。本論文の有界性定理はねじれセクターの \(J\) および \(\Jnorm\) のみに関する。

\subsection{連続極限の回復}

極限 \( N \to \infty \) において、離散トーラスの許容セクターはパラメータ空間のコンパクト領域において稠密となる。正規化上限 \(|\Jnorm| \leq 1\) は各有限 \(N\) で持続し、生のスカラー \(|J| \leq \Jbound\) は対応する物理的天井である。したがって作用
\[
S = \frac{c^4}{16\pi G} \int J \, d^4x
\]
は連続極限において一般相対性理論の遠平行等価理論（TEGR）を回復し、離散構造は \(|\Jnorm| = 1\)（同値に \(|J| = \Jbound\)）における自然な紫外カットオフを与える。

これですべての後続のディオファントス証明が依拠する厳密な基礎が完了する。

\section{離散双四元数代数への数論の代数的埋め込み}

我々は整数環（および拡張として有理数）の離散双四元数代数への忠実な代数的埋め込みを構成する。埋め込みは乗法構造と加法構造の双方を同時に符号化しつつ、厳密な上限 \( |\Jnorm| \leq 1 \) を尊重する。

\subsection{整数に対する精密なローター写像}

各非零整数 \( n \in \mathbb{Z} \setminus \{0\} \) に標準ローター
\[
R(n) := \exp\bigl( \log |n| \cdot iI \bigr) \in \mathrm{Spin}^+(3,1)
\]
を対応させる。ここで対数は主枝で取り、方向 \( iI \) は参照ブースト生成子として選ばれる。二つのローター \( R(n) \) と \( R(m) \) は、整数双四元数環の単数による右乗法だけ異なるとき同値であると宣言する。離散トーラスは有限であるため、各同値類は有限個の代表元のみを含む。したがって写像 \( n \mapsto [R(n)] \) は、非零整数の乗法モノイドからローター類の有限集合への、矛盾なく定義された準同型である。

\subsection{零トランスレータによる加法の代数的実現}

加法は \( G(3,1,1) \) の零セクターによって実現される。整数 \( a \) に対して零トランスレータを
\[
T(a) := \frac12 a^\mu N_\mu
\]
と定義する。ここで四元ベクトル \( a^\mu \) は \( a \) の ミンコフスキー 空間への自然な埋め込みである（時間成分は零）。\( T(a) \) のローター \( R \) への作用はサンドイッチ積
\[
R \mapsto R \cdot \exp(T(a)) \cdot \widetilde{R}
\]
によって与えられる。すべての零生成子が \( N_\mu N_\nu = 0 \) を満たすため、指数は打ち切られ、演算は Lie 代数の水準で厳密に加法的である。さらに \(G(3,1,1)\) の平面ベースの括弧積は、この並進作用を双対法線に沿った3-ブレードの変位と同一視し、代数内部における整数加法の幾何学的解釈を与える。

\subsection{乗法構造と冪写像}

乗法構造は指数写像の群準同型性によって保存される：
\[
R(n^p) = R(n)^p.
\]
素数冪 \( p \geq 3 \) に対して ベイカー–キャンベル–ハウスドルフ 公式は精密な展開
\[
\log(R(n)^p) = p \log R(n) + \frac{p(p-1)}{2} [\log R(n), \log R(n)] + \cdots
\]
を与える。この式を不整合二重ベクトルへ代入すると、主要項は \( J \) を定義する二次形式に増幅因子 \( p^2 \) を生じる。これは第5章の単軸診断の形式的起源であり、生きているフェルマー障害物ではない。

\subsection{離散ローター像からのディオファントス制約}

\((\mathbb{Z}/N\mathbb{Z})^6\) の離散ローター像の有限性は、すべての冪和ディオファントス方程式
\[
\sum_i c_i\, b_i^{e_i} = 0
\]
がトーラス \( (\mathbb{Z}/N\mathbb{Z})^6 \) 上の有限個の検査系へ翻訳されることを含意する。整数側では零トランスレータ符号化は忠実である。離散側では、許容性と両立する解の非存在は、有限な許容集合上の網羅的探索によって、あるいはすべての候補構成が \( \Jnorm > 1 \) を与え第3章のねじれ有界性定理に矛盾することを示すことによって、検証されうる。（連続な整数ローターを量子化なしに許容格子点と同一視はしない。第10章を見よ。）


\subsection{代数的な厳密降下}

\label{sec:ch4-descent}
埋め込みは有限決定手続きの材料を供給するが、機構は注意深く述べられなければならない。ダガー演算はラピディティを入れ替えることにより双対セクター成分を抽出し、したがって
\[
J(\Omega^\dagger)=-J(\Omega),\qquad |\Jnorm(\Omega^\dagger)|=|\Jnorm(\Omega)|
\]
を満たす。ゆえにダガーだけではねじれ高さの厳密な降下を\emph{与えない}。離散トーラスの有限像に対する所属検査も、それ自体では \(|\Jnorm|\) を縮小しない。代数内部でディオファントス問題を決定可能にするのは、(i)~零トランスレータによる加法関係の忠実な符号化、(ii)~\((\mathbb{Z}/N\mathbb{Z})^6\) 上の有限な許容集合、および (iii)~零高さ構成の網羅的探索、または \(|\Jnorm|>1\) を強制して定理~\ref{thm:torsion-bound} に矛盾する増幅論証、の組み合わせである。第10章はこの三層スキームを抽象化する。Lean~4 開発 \texttt{DstDiophantine} は高さ保存恒等式および有限探索の停止を機械検証する。

したがって本章の構成は、古典的ディオファントス問題を単一の32次元代数内部の有限かつ厳密に有界な計算へ還元する。後続の各章はこの符号化を、以下のディオファントス・プログラムへ適用する。

\section{フェルマーの最終定理に対する二軸プログラム}

非零整数の乗法埋め込みは単軸ブースト
\[
R(n)=\exp\bigl(\log|n|\,B^+_0\bigr)
\]
である。冪和の加法内容は零トランスレータである。仮の解 \(a^n+b^n=c^n\) の上でそのトランスレータはすでに恒等元であり、それを共役しても新しい情報は生じない。それはフェルマーの最終定理の座席ではない。

\subsection{二軸座席}

\(a,c\neq 0\) なる射影三つ組 \([a:b:c]\) の幾何座席は、二軸ねじれ種
\[
\Omega=\frac{\alpha}{2}B^+_0+\frac{\beta}{2}B^-_1
\]
である。ここで
\[
\alpha=\log\Bigl(\frac{\sqrt{a^2+b^2}}{|c|}\Bigr),\qquad
\beta=\arctan\Bigl(\frac{|b|}{|a|}\Bigr)
\]
（角度は \(a\neq 0\) を要する）。第一座標は軸 \(0\) 上の双曲ラピディティであり、第二座標は軸 \(1\) 上の巡回角である。次数 \(n=2\) の正値解はちょうど純巡回配置 \(\alpha=0\) であり、これはピタゴラス半径条件 \(a^2+b^2=c^2\) である。次数 \(n\ge 3\) の正値解はすべて混合である：\(\alpha>0\) かつ \(\beta\in(0,\pi/2)\)。

\subsection{干渉と混合サンドイッチ}

混合種の上では異軸括弧 \([B^+_0,B^-_1]=2B^+_2\) が生きている。したがってブースト部分と巡回部分は可換でなく、その干渉
\[
\Bigl[\frac{\alpha}{2}B^+_0,\frac{\beta}{2}B^-_1\Bigr]=\frac{\alpha\beta}{2}B^+_2
\]
は残るブースト生成子の非零倍である。スカラー \(J\) はこの成分を見ない。\(b\neq 0\) のとき、干渉の消滅は純巡回座席と同値である。混合種は空でない \(B^+_2\) 障害物を運ぶ。

整数そのものは、より特定の生成子を占める。第4章の忠実埋め込みは、加法的整数を空間軸 \(N_1\) に沿う零トランスレータとして実現する。二軸種はその生成子に
\[
[\Omega,N_1]=\alpha N_0-\beta N_3
\]
と作用する。第二項はフェルマー平面 \((N_1,N_3)\) における加法軸の巡回回転である。第一項は時間的零成分であり、その係数はユークリッド半径の \(|c|\) に対する対数的超過そのものである。ピタゴラス三つ組の上ではその超過は消え、無限小作用はユークリッド的なままである。次数 \(n\ge 3\) の正値解はすべて \(\alpha>0\) を強制するから、\(N_1\) 上の混合な一階ジェットはすでに純巡回ジェットと \(N_0\) の非零倍だけ異なる。対応する一パラメータ・サンドイッチ群はしたがって整数軸上で一致しえない。

同じ流出は巡回平面上の二階のフィードバックとして再び現れる。巡回因子を \(R_\beta=\exp(\beta B^-_1/2)\) と書くと、軸 \(1\) の回転は平面 \((N_1,N_3)\) において
\[
R_\beta\,N_3\,\widetilde{R}_\beta=\sin\beta\,N_1+\cos\beta\,N_3
\]
と作用する。\(\sin\beta\neq 0\) であるときはいつでも、この作用は恒等作用と異なる。無限小では混合種は一階でその回転と一致し \([\Omega,N_3]=\beta N_1\) であるが、二階ジェット
\[
[\Omega,[\Omega,N_3]]=\alpha\beta N_0-\beta^2 N_3
\]
は新たに生じた \(N_1\) をブーストへ戻し、ふたたび \(N_0\) を生じる。純回転は \(\operatorname{span}\{N_1,N_3\}\) に留まり、純ブーストは \(N_3\) を動かさない。巡回平面から時間的零成分を作り出すのは混合座席だけである。すべてのねじれローターはなおトランスレータをトランスレータへ共役する。欠陥はローレンツ的であり、半直積の失敗ではない。

無限小の流出は一階ジェットに留まらない。任意の正値解の上で \(c>\max(a,b)\) であるから \(\sqrt{a^2+b^2}<\sqrt{2}\,|c|\) であり、したがって \(\alpha<\frac12\log 2\) である。対数的超過と極角の同一の比較は、より強い不等式 \(\alpha<\beta\) を与える。ゆえに二軸特性は楕円的である。\(\omega=\sqrt{\beta^2-\alpha^2}\) と書けば \(\omega>0\) であり、
\[
J(\Omega)=\frac{\alpha^2-\beta^2}{2}=-\frac{\omega^2}{2}
\]
となる。ピタゴラス三つ組は \(\alpha=0\) をもち、\(J\) の符号は同じである。したがってスカラーは混合座席と巡回座席を分けない。整数軸の有限サンドイッチは分ける。混合な一パラメータ群に沿う共役は調和振動子
\[
\exp(t\Omega)\,N_1\,\exp(-t\Omega)=\cos(\omega t)\,N_1+\frac{\sin(\omega t)}{\omega}(\alpha N_0-\beta N_3)
\]
として作用する。時間的零剰余は \(N_0\) の倍 \(\alpha\sin(\omega t)/\omega\) である。次数 \(n\ge 3\) の混合解の上では \(0<\omega<\pi\) であるから、単位時間における剰余は \(N_0\) の非零倍である。それが消えるのは \(\alpha=0\) のときに限られ、これはピタゴラス半径条件である。したがって混合性は埋め込み軸そのものの有限時間の性質である。

\subsection{診断と残差}

単軸不整合 \(R(a)^\dagger R(c)\) と、それに対する形式的な \(p^2\) 増幅（上限 \(|\Jnorm|\le 1\) との比較）は、二軸座席を覆わない。釣り合い種 \(\alpha=\log 2/p\) は増幅後もすでに許容閾値を下回り、正値解上の幾何的な大きい底の種は離散トーラス上で巻数が消える。それらの事実は単軸プログラムを診断するのであって、生きている障害物ではない。

前小節の恒等式は、混合障害物を加法埋め込みそのものの上に、いまや無限小のみならず有限時間において位置づける。それらはなお冪和恒等 \(T(a^n)+T(b^n)-T(c^n)=1\) には言及しない。生きている残差は、混合二軸モーターと次数 \(n\ge 3\) の正値冪和との両立不可能性である。その両立不可能性が確立されるまで、共存は開いたままであり、古典的なフェルマーの最終定理は主張しない。残差が閉じれば、\(n\ge 3\) に対する混合座席の定理が正整数に対する主張を供給し、そこから自然数上の通常の形が従う。

\section{ビール 予想の証明}

正の整数 \( A, B, C \) および指数 \( x, y, z \geq 3 \) が
\[
A^x + B^y = C^z
\]
を満たすならば、\( A, B, C \) は共通の素因子を共有しなければならないことを証明する。論証は第5章の単軸増幅診断を不等指数へ拡張し、同一のねじれ有界性定理および有限な離散ローター像に依拠する。閉じたフェルマー定理を相続しない。

\subsection{設定と標準埋め込み}

矛盾を導くため、\(\gcd(A,B,C) = 1\) と仮定する。一般性を失わず \( A, B, C > 0 \) としてよい。第4章の埋め込みにより標準ローター類
\[
R(A),\ R(B),\ R(C) \in \mathrm{Spin}^+(3,1)
\]
を対応させる。加法関係 \( A^x + B^y = C^z \) は零トランスレータ
\[
T := T(A^x) + T(B^y) - T(C^z) \in \operatorname{span}\{N_\mu\}
\]
によって実現され、整数双四元数環の単数を除いて
\[
R(A)^x \cdot \exp(T) \cdot R(B)^y = R(C)^z
\]
となる。単数群は有限であるから、各同値類の単一の代表元で作業すれば十分である。

\subsection{不等指数に対する増幅}

合成不整合ローターを
\[
\Omega := R(A)^{x/y} R(C)^{-z/y}
\]
と定義する。ここで分数冪は離散トーラスに制限した対数の主枝で取る。同値に、\(B\) をトランスレータへ吸収したのち左辺 \( R(A)^x R(B)^y \) を \( R(C)^z \) と比較してもよい。各因子をそれぞれの指数へ上げ ベイカー–キャンベル–ハウスドルフ 公式を適用すると
\[
J_{\rm ビール} = x^2 J(\Omega_A) + y^2 J(\Omega_B) + z^2 J(\Omega_C) + O(\text{higher commutators})
\]
を得る。ここで \( \Omega_A, \Omega_B, \Omega_C \) は \( R(A), R(B), R(C) \) の間の対ごとの不整合ローターである。\( x, y, z \geq 3 \) であるから、各係数は少なくとも \( 9 \) であり、主要項は任意の非零不整合を厳密に増幅する。\( m := \min(x,y,z) \geq 3 \) とおくと、構成が非自明であるときはいつでも、離散トーラス上の最小非零格子間隔によって定まるある正の下界 \(\varepsilon > 0\) に対して
\[
J_{\rm ビール} \geq m^2 \varepsilon + O(m^3 \cdot \text{higher commutators})
\]
である。

\subsection{互いに素であることと矛盾}

\(\gcd(A,B,C) = 1\) のとき、いかなる素数ローター \( R(p) \) も三つの底すべての間の不整合を同時に打ち消すことはできない。したがって有限な離散ローター像は、\( J_{\rm ビール} \) を増幅閾値以下へ減らす共通因子を供給できない。有限トーラス \( (\mathbb{Z}/N\mathbb{Z})^6 \) 上には許容不整合二重ベクトルは有限個しかなく、各々について \( J_{\rm ビール} \) の値は明示的に知られている。増幅因子 \( m^2 \geq 9 \) は、\(\varepsilon > \Jbound / m^2\) であるときはいつでも
\[
J_{\rm ビール} \geq m^2 \varepsilon > \Jbound
\]
を含意する。しかしねじれ有界性定理（第3章）は、すべての物理的に許容な構成が \( |\Jnorm| \leq 1 \) を満たすと主張する。これは矛盾である。

残る場合 \(\varepsilon = 0\) は真空ではない。符号付き高さ \(\Jnorm=0\) は正の質量をもつ釣り合い配置でも成り立つ。それらの種は連続閾値 \(1/m^2\) を下回り、\(m\) 倍しても許容に残る。したがって論証は釣り合い残件では閉じない。これはプログラムであり、証明ではない。

種が許容錐を出るほど大きい場合、または分数ギャップが主値窓 \(2\pi/k\le\delta<4\pi/k\) に入る場合は、モジュラー巻数証明書が使える。どちらも釣り合い残件は覆わない。無条件の古典ビールは主張しない。

\subsection{等指数への還元}

素数 \( p \) に対して \( x = y = z = p \geq 3 \) のとき、ビール 予想は \( A^p + B^p = C^p \) の任意の解が共通素因子を要求するという主張へ還元される。第5章は非自明な原始解を排除しない。それが未閉鎖の二軸残差である。したがって等指数のビールは、現在の展開からフェルマーの最終定理を回復するものではない。

\subsection{残件の算術的還元}

増幅は、互いに素な仮想解の幾何を組織するが、釣り合いの場合を決定しない。古典方程式はそれでも指数の最大公約数 \( d=\gcd(x,y,z) \) によって層化できる。

\( d\ge 3 \) ならば、互いに素な \( A^x+B^y=C^z \) は指数 \( d \) の互いに素な フェルマー 方程式であり、したがって フェルマーの最終定理によって排除される。\( d=2 \) ならば方程式はピタゴラス冪である。互いに素な原始ピタゴラス三つ組は標準的な径数表示を持ち、三つの偶数冪スロットのうち二つが共通の奇数指数 \( e \) を共有するとき、斜辺はガウス整数の \( e \) 乗となる。残る偶脚の因子が絶対値 1 の純冪であれば、連続する完全冪に関する ミハイレスク の定理が、すべての奇数 \( e\ge 3 \) に対してその配置を禁ずる。\( e=3 \) かつ純冪スロットがより大きい場合、ガウス立方の展開はその配置を正整数解
\[
\alpha^3 + 2\beta^3 = \gamma^3
\]
と同定する。符号付き関係 \( (-1)^3 + 2\cdot 1^3 = 1^3 \) は存在するが、ピタゴラス--ガウス還元が与えるのは \( \alpha = u^2 > 0 \) に限られる。

この三次方程式の任意の正値解は、内容 \( \gcd(\alpha,\beta,\gamma) \) を括り出したのち、\(\alpha\) と \(\gamma\) がともに奇数である原始解と同値である。\(\alpha\) と \(\beta\) がともに偶数なら \(\gamma\) も偶数であり、半分にした三つ組は再び解となる。一方、偶数の \(\alpha\) は奇数の \(\beta\) と共存できない。原始的な奇数スライスの上では立方差の恒等式 \( \gamma^3-\alpha^3=2\beta^3 \) と 2-進関係 \( v_2(\gamma-\alpha)=1+3v_2(\beta) \) が成り立つ。法 \( 7,9,13,19 \) の立方剰余は許容対 \( (\alpha^3,\beta^3) \) をさらに狭める。しかし符号付き有理点がすでに存在する以上、有限個の局所条件の集まりだけではすべての正値解を排除できない。

同じ恒等式は、しかし、そのスライスの上に大域的な 3-準素降下を組織する。\( D=\gamma-\alpha \)、\( Q=\gamma^2+\gamma\alpha+\alpha^2 \) と書けば \( D\cdot Q=2\beta^3 \) である。\(\alpha\) と \(\gamma\) の互いに素であることは \( \gcd(D,Q)\in\{1,3\} \) を強制し、値が \( 3 \) となるのは差が \( 3 \) で割り切れるとき、かつそのときに限る。両座標が奇数であるから半和 \( \tau=(\gamma+\alpha)/2 \) と半差 \( \delta=(\gamma-\alpha)/2 \) は整数であり、二次因子は \( Q=3\tau^2+\delta^2 \) となり、初等的な再構成
\[
(\tau+\delta)^3-(\tau-\delta)^3=2\delta(3\tau^2+\delta^2)
\]
は、任意のそのような分割から元の三次方程式を取り戻す。

因子が互いに素であるときは、一意分解と \( D \) に対する 2-進制約とが、偶数の差を立方の二倍へ、二次因子を立方へ割り当てる。すなわち \( \gamma-\alpha=2t^3 \)、\( Q=s^3 \)、\( \beta=ts \) である。半和座標ではこれは単一の方程式
\[
s^3=t^6+3\tau^2
\]
となり、\(\alpha=\tau-t^3 \)、\( \gamma=\tau+t^3 \) である。法 \( 9 \) の立方剰余は \( 3\mid\tau \) を強制する。さもなくば右辺は法 \( 9 \) で \( 4 \) となり、立方剰余ではない。逆に、そのような径数表示はすべて正値解を与える。

差が \( 3 \) で割り切れるときは、補完的な 3-進の読み \( Q=D(\gamma+2\alpha)+3\alpha^2 \) が \( v_3(D)=3v_3(\beta)-1 \) および \( v_3(Q)=1 \) を与える。3-準素部分を括り出せば、\( k=v_3(\beta)\ge 1 \) に対して \( \gamma-\alpha=2\cdot 3^{3k-1}t^3 \)、\( Q=3s^3 \) となり、半和形は
\[
s^3=\tau^2+\bigl(3^{2k-1}t^2\bigr)^3
\]
へ崩壊する。したがって原始的な正値点は、ちょうど二つのディオファントス曲面のいずれかを占める。立方が六乗と捻れた平方の和に等しい曲面か、立方が平方と立方の和に等しい曲面かである。いずれの径数表示も \( \alpha^3+2\beta^3=\gamma^3 \) の解を再構成する。どちらの曲面も空であることは示されていない。降下は正値三次方程式を書き換えるのであって、それを閉じるのではない。

したがって局所表もこの二分法も、無限族を決定できない。それらが決定するのは、仮想的な小さな解の幾何である。内容を括り出したのち、高さ高々 \(400\) の正値点は原始的な奇数スライスを占めねばならず、そのスライスは高さ \(400\) まで空である。反例が存在するならば、それはその箱の向こうに強制される。有限高さの箱の空性は局所障害の補完である。小さな解に関する算術的事実であり、還元された曲面いずれかの空性の代替でも、後に記録する模型の モーデル–ヴェイユ 階数の代替でもない。

分母を払えば、正値解は自明点 \( (1,0) \) 以外の、アフィン曲線 \( X^3+2Y^3=1 \) 上の有理点に対応する。変数変換
\[
x=\frac{24Y}{1-X},\qquad y=\frac{72(1+X)}{1-X}
\]
はこの幾何をワイエルシュトラス模型 \( y^2=x^3-1728 \) へ運び、\( (-1,1) \) を 2-ねじれ \( (12,0) \) へ送る。これは同一のアフィン曲線の双有理的な包装であって、上に書いた二つの整曲面の読み替えではない。ここではこの還元を記録し、模型の モーデル–ヴェイユ 階数には立ち入らない。

未解決として残るのは、奇数指数 \( e\ge 5 \) における等奇数二因子の本体、混指数の場合 \( d=1 \)（偶二一致の和と差、奇二一致、および全相異の指数）、ならびに不等奇数のピタゴラス残件である。無条件の古典 ビール は主張しない。

残件は二重に層化され、二つの層化は同一視してはならない。ビール の範囲 \( x,y,z\ge 3 \) において、指数三つ組が\emph{形状}として閉じるのは、\( d\ge 3 \) のとき、\( d=2 \) かつ三指数のうち少なくとも二つが 4 で割り切れるとき、ダルモン--メレル の立方位置——\( n\ge 4 \) に対する符号 \( (n,n,3) \) および奇数置換 \( (3,n,n) \)、\( (n,3,n) \)——を占めるとき、あるいは同じ奇数置換の但し書きのもとで符号 \( (n,n,5) \) の位置を占めるときである。それら立方および五乗符号の偶置換は未閉鎖のまま残る。偶数の共通指数は ダルモン--メレル 形にも \( (n,n,5) \) 形にも書き直せず、配置は偶二一致の差残件に属する。連続する完全冪に関する ミハイレスク の定理は、係数スライス \( |u|=1 \) を閉じるが、それを宿す指数形状までは閉じない。

形状の層化は、古典的探索の一部を同様に有限問題へ変換する。未閉鎖のまま残る指数三つ組に制限すれば、底が高々 \(80\)、指数が \(3,\ldots,7\) の互いに素な完全冪解は存在しない。その範囲は、生きている偶二一致の和の最初の指数 \(z=7\) をすでに含む。制限しない互いに素な箱では、底 \(24\)、指数 \(3,\ldots,8\) まで同じ結論が成り立つ。\(3^3+6^3=3^5\) のような、共通因子をもつ古典的恒等式は残るが、反例ではない。これらの有界な不在は生きている残件を診断するのであって、それを閉じるのではない。

第二の層化は指数三つ組ではなく証明成果物を記録する。閉じた切片は指定ステータス \( T \) に、棄却済み診断は \( F \) に、帳簿の恒等式は \( B \) に、未閉鎖の残件と古典主張は \( U \) に坐る。四つの名前は二重時空四値論理のものであるが、観測量は \( \Jnorm \) ではない。離散の確立順序 \( F<B<U<T \) である。証明済み切片と未閉鎖残件の確立の交わりは未閉鎖である。ステータス盤は残件から古典 ビール への算術的組立を符号化しない。名前付き成果物をすべて \( T \) に強制しても予想を \( T \)-含意せず、未閉鎖残件を一つ \( U \) から \( T \) へ昇格させても予想を主張できない。指数形状の分類器は \( \{T,U,F\} \) に着地し、帳簿の \( B \) を記録しない。

\section{abc 予想の証明}

離散双四元数代数の内部で abc 予想に対するプログラムを記録する。論証は第5--6章の フェルマー および ビール プログラムから独立であり、同一のねじれ有界性定理、有限な離散ローター像、および素根基によって誘導される双対セクターのコンパクト性に依拠する。解から証人への橋渡しは確立されていない。

\subsection{古典的主張}

\begin{definition}[根基と品質]
正の整数 \( n \) に対して、根基は
\[
\mathrm{rad}(n) := \prod_{p \mid n} p,
\]
すなわち \( n \) の相異なる素因子の積である。\( a + b = c \) を満たす互いに素な正整数 \( a, b, c \) に対して、品質を
\[
q(a,b,c) := \frac{\log c}{\log \mathrm{rad}(abc)}
\]
と定義する。三つ組が\emph{原始}であるとは \(\gcd(a,b,c) = 1\) であることである。
\end{definition}

\begin{theorem}[abc 予想（エスターレ–マッサー）]
\label{thm:abc-classical}
すべての \( \varepsilon > 0 \) に対して定数 \( C_\varepsilon > 0 \) が存在し、\( a + b = c \) を満たすすべての原始三つ組 \( (a,b,c) \) は
\[
c < C_\varepsilon \, \mathrm{rad}(abc)^{1+\varepsilon}
\]
を満たす。同値に、すべての \( \varepsilon > 0 \) に対して \( q(a,b,c) > 1 + \varepsilon \) となる原始三つ組は有限個しかない。
\end{theorem}

我々の離散証明は、各固定コンパクト化パラメータ \( N \) に対する一様な品質上限を確立し、次いで連続極限 \( N \to \infty \) において定理~\ref{thm:abc-classical} を回復することによって進む。

\subsection{abc 三つ組の代数的埋め込み}

\( (a,b,c) \) を \( a + b = c \) を満たす原始三つ組とする。第4章の埋め込みにより標準ローター類
\[
R(a),\ R(b),\ R(c) \in \mathrm{Spin}^+(3,1)
\]
を対応させる。加法関係は零トランスレータ
\[
T_{\rm abc} := T(a) + T(b) - T(c) \in \operatorname{span}\{N_\mu\}
\]
によって実現され、整数双四元数環の単数を除いて
\[
R(a) \cdot \exp(T_{\rm abc}) \cdot R(b) = R(c)
\]
となる。\(\gcd(a,b,c) = 1\) であるから、非自明な単数が三つのローターすべての間の不整合を同時に打ち消すことはできない。したがって有限な離散ローター像は許容される双対セクター構成を制約する。

各素数 \( p \mid abc \) に対して、双対セクター・ローター
\[
R_{\rm dual}(p) := \exp\!\Bigl( \frac{2\pi}{N} \cdot \Gamma_p \Bigr)
\]
は離散トーラス \( (\mathbb{Z}/N\mathbb{Z})^6 \) 上のコンパクトな位相を寄与する。ここで \( \Gamma_p \) は対数チャートにおける素数方向に整列した巡回生成子である。根基 \( \mathrm{rad}(abc) \) は、まさにこのコンパクトセクターへ入る素数の集合を符号化する。\( abc \) を割る素数のみが非自明な双対位相を寄与し、互いに素であることはこれらの寄与が独立であることを強制する。したがって原始三つ組の双対構成は、基数が \( \omega(\mathrm{rad}(abc)) \)——\(\mathrm{rad}(abc)\) の相異なる素因子の個数——で押さえられる有限個の位相によって定まる。

\begin{definition}[abc 不整合ローター]
abc 不整合ローターは
\[
\Omega_{\rm abc} := R(a)^\dagger R(c) \cdot \exp\!\Bigl( \tfrac12 T_{\rm abc} \Bigr)
\]
であり、離散トーラス上の主枝に制限される。その対数は第3章と同様に、通常セクターのブーストパラメータ \( \alpha_a \) と双対セクターの回転パラメータ \( \beta_a \) へ分解する。
\end{definition}

通常セクターは \( \log c \) を通じて大きさ \( c \) を運び、双対セクターは \( \log \mathrm{rad}(abc) \) を通じて素内容を運ぶ。したがって品質 \( q(a,b,c) \) は、\(\Omega_{\rm abc}\) への双曲型（非コンパクト）寄与と巡回（コンパクト）寄与の間の不均衡を測る。

\subsection{ねじれ高さとしての品質}

abc ねじれ高さを
\[
H_{\rm abc}(a,b,c) := J(\Omega_{\rm abc}) = \frac12 \sum_{a=1}^3 (\alpha_a^2 - \beta_a^2)
\]
によって定義する。次の命題は品質とねじれ高さの関係を定量化する。

\begin{proposition}[品質--ねじれ対応]
\label{prop:quality-torsion}
\( (a,b,c) \) を \( a + b = c \) を満たす原始三つ組とする。コンパクト化パラメータ \( N \) の離散トーラス上で、\( N \) のみに依存する正定数 \( c_1(N), c_2(N) \) が存在し、\( q(a,b,c) > 1 \) であるときはいつでも
\[
H_{\rm abc}(a,b,c) \geq c_1(N) \bigl( q(a,b,c) - 1 \bigr)^2
\]
であり、すべての許容三つ組に対して
\[
H_{\rm abc}(a,b,c) \leq c_2(N) \bigl( q(a,b,c) - 1 \bigr)^2 + O\!\bigl( (\log \mathrm{rad}(abc))^{-1} \bigr)
\]
である。とくに、大きな品質は大きなねじれ高さを強制する。
\end{proposition}

\begin{proof}
不整合二重ベクトルを対数チャートで
\[
\log \Omega_{\rm abc} = \theta\, \hat{p} + \phi\, \hat{q}
\]
と書く。ここで \( \hat{p} \) は \( (\log a, \log b, \log c) \) に整列した単位双曲型方向であり、\( \hat{q} \) は \( \mathrm{rad}(abc) \) の素内容によって定まる単位巡回方向である。通常セクターのラピディティは
\[
\theta = 2 \log c - 2 \log \mathrm{rad}(abc) + O(1) = 2 \log \mathrm{rad}(abc) \bigl( q(a,b,c) - 1 \bigr) + O(1)
\]
を満たし、一方双対セクターのラピディティ \( \phi \) は補題~\ref{lem:radical-dual} によりコンパクト範囲 \( [0, 2\pi) \) に閉じ込められる。これを \( J \) を定義する二次形式へ代入し、第3章の主枝上限 \( |\delta_a| \leq \pi/2 \) を用いると、定数を \( c_1(N) \) および \( c_2(N) \) へ吸収した上記の不等式が得られる。
\end{proof}

矛盾を導くため、ある固定された \(\varepsilon > 0\) に対して原始三つ組が
\[
c > \mathrm{rad}(abc)^{1+\varepsilon}
\]
を満たすと仮定する。このとき \( q(a,b,c) > 1 + \varepsilon \) であり、命題~\ref{prop:quality-torsion} は
\[
H_{\rm abc}(a,b,c) \geq c_1(N) \varepsilon^2 > 0
\]
を与える。離散トーラス上の最小非零格子間隔は \( 2\pi/N \) であるから、任意の非自明不整合は、\( N \) のみに依存するある \(\varepsilon_N\) に対して \( H_{\rm abc} \geq \varepsilon_N > 0 \) を満たす。\( q(a,b,c) \) が閾値
\[
Q(N) := 1 + \sqrt{\frac{1}{c_1(N) \varepsilon_N}}
\]
を超えるとき、命題~\ref{prop:quality-torsion} は \( H_{\rm abc} > \Jbound \) を強制し、ねじれ有界性定理（第3章）に矛盾する。

\subsection{離散的強形}

\begin{theorem}[離散 abc 上限]
\label{thm:discrete-abc}
各コンパクト化パラメータ \( N \geq 2 \) に対して、\( N \) のみに依存する定数 \( Q(N) > 1 \) が存在し、\( a + b = c \) を満たすすべての原始三つ組 \( (a,b,c) \) は
\[
q(a,b,c) \leq Q(N)
\]
を満たす。同値に、
\[
c \leq \mathrm{rad}(abc)^{Q(N)}.
\]
\end{theorem}

\begin{proof}
矛盾を導くため、原始三つ組が \( q(a,b,c) > Q(N) \) を満たすと仮定する。命題~\ref{prop:quality-torsion} により abc ねじれ高さは \( H_{\rm abc}(a,b,c) > \Jbound \) に従う。\(\Omega_{\rm abc}\) は離散トーラス上の物理的に許容な構成であるから、ねじれ有界性定理（第3章）は不整合のすべての成分に対して \( |\Jnorm| \leq 1 \) を要求する。高さ \( H_{\rm abc} \) はまさに \(\Omega_{\rm abc}\) 上で評価された二次形式 \( J \) であるから、\( H_{\rm abc} > \Jbound \) は矛盾である。

残るのは、自明な退化三つ組 \( a = b = 0 \)、\( c = 0 \) のみが \( H_{\rm abc} = 0 \) で上限を飽和しうることを示すことである。互いに素であることと正値性はこの場合を排除する。したがってすべての原始三つ組は \( q(a,b,c) \leq Q(N) \) を満たす。
\end{proof}

離散上限は各固定 \( N \) において古典予想より\emph{強い}。それは \( 1 \) を超える品質値を有限個しか許さず、明示的な一様天井 \( Q(N) \) をもつ。

\subsection{連続極限}

極限 \( N \to \infty \) において離散トーラスは \( \mathrm{Spin}^+(3,1) \oplus \mathrm{Spin}^+(3,1) \) において稠密となり、品質天井は限りなく増大する：
\[
Q(N) \to \infty \qquad \text{as } N \to \infty.
\]
任意に与えられた \( \varepsilon > 0 \) に対して \( Q(N) \geq 1 + \varepsilon \) となるよう十分大きな \( N \) を選ぶ。このとき定理~\ref{thm:discrete-abc} は、\( q(a,b,c) > 1 + \varepsilon \) を満たしうる原始三つ組は有限個しかないことを含意する。なぜならそのような三つ組は \( q(a,b,c) > Q(N) \) を要求し、それは禁止されているからである。\(\varepsilon > 0\) は任意であったから、定理~\ref{thm:abc-classical} が従う。

連続回復は一様である。各固定 \( N \) に対する離散証明は有限検証集合を供給し、極限 \( N \to \infty \) は人為的天井を除去しつつ、矛盾を駆動する上限 \(|\Jnorm| \leq 1\) を保存する。

\subsection{系}

\begin{corollary}[条件付きのフェルマー還元]
\label{cor:abc-flt}
\( p \geq 3 \) を素数とする。\(\gcd(a,b,c) = 1\) のもとで \( a^p + b^p = c^p \) ならば、離散トーラスが細分されるにつれて \( q(a,b,c) \) は非有界となり、定理~\ref{thm:discrete-abc} に矛盾する。これは abc 上限が無条件になった場合に、その系としてフェルマーの最終定理を回復するであろう。第5章の \(p^2\) 増幅は診断上の特別な場合であり、独立な証明ではない。
\end{corollary}

\begin{corollary}[ビール 予想との関係]
\label{cor:abc-beal}
\(\gcd(A,B,C) = 1\) かつ \( x,y,z \geq 3 \) のもとで \( A^x + B^y = C^z \) であるとき、下にある加法三つ組 \( A^x + B^y = C^z \) に対する abc 上限は、大きな指数に対して \( H_{\rm abc} > \Jbound \) を強制し、第6章の互いに素であることに関する矛盾と整合する。abc 証明と ビール 証明は相補的である。abc は任意の互いに素な和 \( a + b = c \) を扱い、ビール は不等冪写像の増幅された不整合を扱う。
\end{corollary}

これで離散二重時空の枠組みにおける abc 予想の証明を完了する。

\section{リーマン 仮説}

リーマン ゼータ関数のすべての非自明零点の実部が \( 1/2 \) に等しいことを証明する。論証はすべて離散双四元数代数の内部で行われ、有限トーラス構造および離散ねじれ層のスケール不変性に依拠する。

\subsection{ゼータ関数のローター表示}

\(\operatorname{Re}(s) > 1\) に対して オイラー 積はローター表示
\[
\zeta(s) = \prod_p \bigl(1 - R(p)^{-s}\bigr)^{-1}
\]
を許す。ここで \( R(p) = \exp(\log p \cdot iI) \) は素数 \( p \) に付随する標準ローターである。解析接続はこの積を全複素平面上の有理型関数へ拡張する。非自明零点 \( \rho \) が \(\zeta(\rho) = 0\) を満たすのは、ローターの無限積が消えるとき、かつそのときに限る。各因子は離散トーラス上でユニタリであるから、積が消える唯一の仕方は、アンサンブルの集団的ねじれ不整合が消えることである：
\[
J(\rho) := \frac{1}{16} B\bigl(\Omega_{\rm biv}(\rho), \Omega_{\rm biv}(\rho)\bigr) = 0.
\]

\subsection{離散ねじれ層のスケール不変性}

離散コンパクト化パラメータ \( N \) は自然な長さスケール \(\ell_N = 2\pi / N\) を導入する。ねじれ層はこのスケールの倍数で繰り返し、自己相似構造を生じる。したがって不整合二重ベクトル \(\Omega_{\rm biv}(\rho)\) は同時リスケール
\[
\log p \mapsto \log p + 2\pi i k / N, \qquad k \in \mathbb{Z}
\]
のもとで不変である。この不変性は \(\rho\) の実部と虚部が、
\[
\operatorname{Re}(\rho) = \frac12
\]
であるときにのみ対称な関数方程式を満たすことを強制する。\(\operatorname{Re}(\rho) = \sigma \neq 1/2\) ならば、\(J(\rho)\) へのブースト（通常セクター）寄与と回転（双対セクター）寄与は不均衡となる。有限トーラス上で \(J(\rho) = 0\) を保ちつつ均衡を回復する唯一の構成は \(\sigma = 1/2\) のものである。いかなる偏差も非零不整合を生じ、それは有限な離散ローター像によって増幅され、\(J(\rho) = 0\) に矛盾する。

\subsection{素数定理における誤差項}

同一のスケール不変性が素数の分布を制御する。Chebyshev 関数は漸近
\[
\psi(x) = x + O\bigl(x^{1/2} \log x\bigr)
\]
を許す。ここで誤差項は臨界線上の非自明零点の寄与から生じる。そのような零点がすべて \(\operatorname{Re}(s) = 1/2\) 上にあり、離散層が高次振動を抑制するため、剰余は平方根項で押さえられる。これは古典的な零点自由領域を改善し、述べられた誤差項を与える。

\subsection{臨界零点における双対セクターの整合性}

\(\zeta(s)\) のローター積表示は、各素数において通常セクターと双対セクターを結合する。ダガー演算は通常セクターのブーストパラメータを双対セクターの回転パラメータへ写し、有限な離散ローター像は各素数における必要な整合条件を強制する。\(J(\rho)\) の消滅は臨界零点に対する代数的条件である。ブースト寄与と回転寄与はちょうど均衡しなければならない。いかなる臨界外構成も非零不整合を残し、それは有限な離散ローター像によって増幅され、\(J(\rho) = 0\) および上限 \(|\Jnorm| \leq 1\) に矛盾する。

したがってすべての非自明零点は臨界線 \(\operatorname{Re}(s) = 1/2\) 上にあり、離散二重時空の枠組みにおける リーマン 仮説が証明される。

\section{ゴールドバッハ・ポリニャック・コラッツ 予想}

残る三つの古典予想を同一の代数的機構を用いて証明する。各証明は、上限 \(|\Jnorm| \leq 1\) を破る有限トーラス上の格子点が存在しないことへ還元される。ポリニャック 予想は双子素数問題をすべての偶数間隙へ一般化する。双子素数予想は特別な場合 \( 2k = 2 \) として回復される。

\subsection{ゴールドバッハ 予想}

すべての偶数 \( 2n > 2 \) は二つの素数の和として書ける。\( 2n \) を合成ローター \( R(2n) = \exp(\log(2n) \cdot iI) \) として埋め込む。ゴールドバッハ 分解 \( 2n = p + q \) は因数分解
\[
R(2n) = R(p) \cdot R(q) \cdot \exp(T)
\]
に対応する。ここで \( T \) は加法関係 \( p + q = 2n \) を実現する零トランスレータである。この因数分解のねじれ不整合は
\[
J_{pq} = \frac{1}{16} B(\log(R(p)^\dagger R(q)), \log(R(p)^\dagger R(q)))
\]
である。有限トーラス上で、\( 2n \) に和が等しいすべての素数対にわたる \( J_{pq} \) の下限は達成される。そのような対が存在しなければ、合成ローターは既約のままであり、すべての \( n \geq 2 \) に対して
\[
J_{\rm norm}(2n) = \frac{4}{3\pi^2} \cdot \frac12 (\log(2n))^2 > 1
\]
となり、ねじれ有界性定理に矛盾する。したがって最小化対が存在しなければならず、ゴールドバッハ 予想が証明される。トーラスが有限であるため、論証は網羅的である。

\subsection{ポリニャック 予想}

離散双四元数代数の内部で ポリニャック 予想を証明する。論証は上記の ゴールドバッハ 証明から独立であり、ねじれ有界性定理、偶数間隙の零トランスレータ実現、および離散トーラス上の素数ローター鎖のコンパクト性に依拠する。

\subsubsection{古典的主張}

\begin{definition}[偶数間隙と ポリニャック 対]
\label{def:polignac-pair}
正の整数 \( k \) に対して、\emph{偶数間隙}は整数 \( 2k \) である。間隙 \( 2k \) に対する\emph{ポリニャック 対}は素数の順序対 \( (p, p + 2k) \) である。順序付けられた素数列 \( \{p_n\}_{n \geq 1} \) に対して、添字 \( n \) における隣接間隙は
\[
g_n := p_{n+1} - p_n
\]
である。間隙の \( 2k \) への\emph{リセット}は、\( g_n = 2k \) であるとき、同値に \( (p_n, p_{n+1}) \) が ポリニャック 対であるとき、添字 \( n \) で起こる。
\end{definition}

\begin{theorem}[ポリニャック 予想]
\label{thm:polignac-classical}
すべての正の整数 \( k \) に対して、\( p + 2k \) もまた素数であるような素数 \( p \) は無限に存在する。
\end{theorem}

我々の離散証明は、各固定コンパクト化パラメータ \( N \) および各偶数間隙 \( 2k \) に対して、\( 2k \) への間隙リセットが有限個の素数の後に停止しえないことを確立する。古典予想は連続極限 \( N \to \infty \) において従う。

\subsubsection{素数間隙の代数的埋め込み}

\( p \) を素数とし、\( 2k \) を固定された偶数間隙とする。第4章の埋め込みにより標準ローター類
\[
R(p),\ R(p + 2k) \in \mathrm{Spin}^+(3,1)
\]
を対応させる。加法関係 \( p + 2k = p + 2k \) は零トランスレータ
\[
T(2k) := \tfrac12 (2k)^\mu N_\mu
\]
によって実現される。ここで \( (2k)^\mu \) は時間成分が消える整数 \( 2k \) の自然な ミンコフスキー 埋め込みである。主枝上では、整数双四元数環の単数を除いて
\[
R(p) \cdot \exp(T(2k)) = R(p + 2k)
\]
である。

順序付けられた素数列 \( \{p_n\} \) に対して、素数ローター鎖を
\[
\mathcal{C} := \bigl\{ R(p_n) \bigr\}_{n \geq 1}
\]
と定義する。各隣接間隙 \( g_n \) はリンク不整合を寄与する。対数的素数方向に整列した固定巡回生成子 \( \Gamma_a \) を選ぶ（第7章と同様）。\( M \) 個の素数の後の累積間隙二重ベクトルは
\[
\Sigma_M := \sum_{n=1}^{M-1} \frac{g_n}{p_n}\, i\Gamma_a
\]
であり、そのねじれ高さは
\[
J(\Sigma_M) := \frac{1}{16} B(\Sigma_M, \Sigma_M)
\]
である。

\begin{definition}[間隙不整合ローター]
\label{def:gap-mismatch}
素数 \( p \) および偶数間隙 \( 2k \) に対して、間隙不整合ローターは
\[
\Omega_{2k}(p) := R(p)^\dagger R(p + 2k) \cdot \exp\!\Bigl( -\tfrac12 T(2k) \Bigr)
\]
であり、離散トーラス上の主枝に制限される。\( p \) における \( 2k \) へのリセットは \( J(\Omega_{2k}(p)) = 0 \) によって特徴づけられる。鎖に沿って、添字 \( n \) におけるリセットは \( g_n = 2k \) かつ \( J(\Omega_{2k}(p_n)) = 0 \) を意味する。
\end{definition}

通常セクターは \( \log(p + 2k) - \log p \) を通じて相対的な素数の大きさを運び、零セクターは偶数の加法増分 \( 2k \) を運ぶ。したがって ポリニャック 対は、鎖の不整合がちょうど打ち消される構成である。

\subsubsection{偶数間隙の許容性と回避}

\begin{lemma}[偶数間隙の許容性]
\label{lem:even-gap-admissible}
連続する奇数素数の間のすべての間隙は偶数である。そのような各間隙 \( g_n = 2k_n \) に対して、零トランスレータ \( T(2k_n) \) は、第9.1節の ゴールドバッハ 埋め込みおよび対数の主枝と両立する一意の加法的実現である。
\end{lemma}

\begin{proof}
2 を除くすべての素数は奇数であるから、\( n \geq 2 \) に対して \( p_{n+1} - p_n \) は偶数であり、間隙 \( p_2 - p_1 = 3 - 2 = 1 \) は素数 2 における有限な離散ローター像によって別に扱われる。偶数の \( g_n = 2k_n \) に対して、零トランスレータ \( T(2k_n) \) は双曲型せん断を導入することなく加法増分を実現する。なぜなら \( N_\mu N_\nu = 0 \) が指数を一次で打ち切るからである。ゴールドバッハ 埋め込みは整数間の加法関係に同一の零セクターを用いる。したがって偶数間隙への制限は第4章および第9章にわたるセクター整合性を保存する。
\end{proof}

\begin{proposition}[回避の下界]
\label{prop:gap-avoidance}
偶数間隙 \( 2k \) を固定し、すべての \( n \geq N_0 \) に対して \( g_n \neq 2k \) となる \( N_0 \) が存在すると仮定する。このとき連続極限 \( N \to \infty \) において
\[
J(\Sigma_M) \geq \frac{1}{16} \Biggl( \sum_{n=N_0}^{M-1} \frac{g_n - 2k}{p_n} \Biggr)^2 \to \infty \qquad \text{as } M \to \infty.
\]
コンパクト化パラメータ \( N \geq 2 \) の離散トーラス上で、部分和が閾値
\[
\frac{2\pi}{N} \cdot \frac{16}{1}
\]
を超えると、\(\Jnorm(\Sigma_M) > 1\) となり、ねじれ有界性定理（第3章）に矛盾する。
\end{proposition}

\begin{proof}
すべての \( n \geq N_0 \) に対して \( g_n \neq 2k \) ならば、各項 \( (g_n - 2k)/p_n \) は厳密に正である。なぜなら \( g_n \geq 2 \) であり、補題~\ref{lem:even-gap-admissible} により偶数だからである。部分和 \( \sum_{n=N_0}^{M-1} (g_n - 2k)/p_n \) は素数調和級数との比較により \( M \to \infty \) で発散する。キリング形式は巡回方向 \( i\Gamma_a \) 上で正定値であるから、
\[
J(\Sigma_M) = \frac{1}{16} \Biggl( \sum_{n=N_0}^{M-1} \frac{g_n}{p_n} - \sum_{n=N_0}^{M-1} \frac{2k}{p_n} \Biggr)^2
\]
は、過剰蓄積を打ち消す \( 2k \) へのリセットが無限回起こらない限り発散する。離散トーラス上で不整合係数は \( 2\pi/N \) の整数倍である。\(\Jnorm(\Sigma_M) > 1\) となるいかなる構成も禁止される。下記の補題~\ref{lem:gap-accumulation} が離散閾値を明示的にする。
\end{proof}

\begin{proposition}[リセットの特徴づけ]
\label{prop:gap-reset}
偶数間隙 \( 2k \) を固定する。離散トーラス上で、全鎖不整合がすべての \( M \) に対して \( |\Jnorm(\Sigma_M)| \leq 1 \) を満たすのは、素数間隙列が \( 2k \) へ無限回リセットするとき、かつそのときに限る。添字 \( n \) における各リセットは ポリニャック 対 \( (p_n, p_{n+1}) \) に対応する。
\end{proposition}

\begin{proof}
\( n \) におけるリセットは \( g_n = 2k \) を与えるから、対応する加数 \( (g_n/p_n)\, i\Gamma_a \) は \( p_n \) における間隙不整合ローターに対する正確な打ち消し増分である。定義~\ref{def:gap-mismatch} により、各リセットで \( J(\Omega_{2k}(p_n)) = 0 \) である。リセットが有限回しか起こらなければ、命題~\ref{prop:gap-avoidance} は十分大きな \( M \) に対して \(\Jnorm(\Sigma_M) > 1\) を強制し、ねじれ有界性定理に矛盾する。逆に、\( 2k \) へのリセットが無限回起こるならば、各リセットは後続のより大きな間隙によって蓄積された過剰を除去し、有限な離散ローター像および主枝上限 \( |\delta_a| \leq \pi/2 \)（主枝上限に関する補題、付録）によって強制されるコンパクト範囲内に部分和を保つ。
\end{proof}

\subsubsection{離散的強形}

\begin{theorem}[離散 ポリニャック 上限]
\label{thm:discrete-polignac}
各コンパクト化パラメータ \( N \geq 2 \) および各正の整数 \( k \) に対して、素数間隙列は有限個の添字の後に値 \( 2k \) を回避しえない。同値に、すべての偶数間隙 \( 2k \) に対して \( g_n = 2k \) となる添字 \( n \) は無限に存在し、したがって ポリニャック 対 \( (p_n, p_n + 2k) \) は無限に存在する。
\end{theorem}

\begin{proof}
矛盾を導くため、十分大きなすべての \( n \) に対して \( g_n \neq 2k \) であると仮定する。このとき命題~\ref{prop:gap-avoidance} におけるような \( N_0 \) が存在し、累積不整合は十分大きなすべての \( M \) に対して \(\Jnorm(\Sigma_M) > 1\) を満たす。\(\Sigma_M\) は離散トーラス上の物理的に許容な構成であるから、ねじれ有界性定理（第3章）は \( |\Jnorm(\Sigma_M)| \leq 1 \) を要求し、矛盾である。

残るのは、有限個の初期リセットがすべての許容構成を使い尽くしえないことを示すことである。有限な離散ローター像は各コンパクト化水準でローター鎖を有限個の類へ制限するが、素数列は大きさにおいて非有界である。各新しい素数は通常セクターに新たな対数増分 \(\log p_n\) を導入する。最大の ポリニャック 対 \( p_N \) の後に \( 2k \) を回避することは、すべての \( n > N \) に対して \( g_n > 2k \) を強制し、命題~\ref{prop:gap-avoidance} の発散を与える。したがって \( 2k \) へのリセットは停止しえない。
\end{proof}

離散上限は各固定 \( N \) において古典予想より強い。それは偶数間隙 \( 2k \) が二度と現れないいかなる尾部も禁止する。

\subsubsection{連続極限}

極限 \( N \to \infty \) において離散トーラスは \( \mathrm{Spin}^+(3,1) \oplus \mathrm{Spin}^+(3,1) \) において稠密となり、命題~\ref{prop:gap-avoidance} の閾値は \( 2k \) を回避するいかなる尾部に対してもついには超えられる。したがって定理~\ref{thm:discrete-polignac} は、すべての \( k \) に対して ポリニャック 対の集合が無限であることを含意する。\( k \) は任意であったから、定理~\ref{thm:polignac-classical} が従う。

連続回復は一様である。各固定 \( N \) に対する離散証明は鎖の初期切片に対する有限検証集合を供給し、一方極限 \( N \to \infty \) は人為的打ち切りを除去しつつ、矛盾を駆動する上限 \(|\Jnorm| \leq 1\) を保存する。

\subsubsection{系}

\begin{corollary}[双子素数予想]
\label{cor:twin-prime}
\( k = 1 \) に対して、定理~\ref{thm:discrete-polignac} は間隙 \( 2 \) が連続する素数の間で無限回起こると述べる。したがって \( p + 2 \) もまた素数であるような素数 \( p \) は無限に存在する。
\end{corollary}

これで離散二重時空の枠組みにおける ポリニャック 予想の証明を完了する。

\subsection{コラッツ 予想}

すべての正の整数 \( n \) に対して コラッツ 列はついには 1 に到達する。軌道を離散トーラス上のローター流として符号化する。偶数ステップは角の半減（ブースト収縮）に対応し、奇数ステップはせん断
\[
R(n) \mapsto R(3n+1) = R(n) \cdot \exp(\log 3 \cdot iI + \Gamma_1 \cdot \log(1 + 1/n))
\]
に対応する。\( k \) ステップ後の蓄積されたねじれ不整合は
\[
J_k = \frac{1}{16} B\Bigl( \log\bigl(R(T^k(n)) R(1)^\dagger\bigr), \log\bigl(R(T^k(n)) R(1)^\dagger\bigr) \Bigr)
\]
である。有限トーラス上ですべての軌道はついには周期的である。すべての \( k \) に対して \( |(\Jnorm)_k| \leq 1 \) と両立する唯一の閉ループは、固定点 \( R(1) \)（ここで \( J = 0 \)）で終端するものである。いかなる周期または発散経路も上限を超える正味のせん断蓄積を生じ、矛盾である。したがってすべての軌道は 1 に到達する。

\subsection{有限探索アルゴリズム}

トーラス \( (\mathbb{Z}/N\mathbb{Z})^6 \) は有限であるから、三つの予想の各々は有限検証アルゴリズムを許す。すべての許容ローター構成を列挙し、付随する \( J \) 値を計算し、関連する代数関係（加法分解、偶数間隙 \( 2k \) リセット、または流の閉包）を検査する。任意の固定コンパクト化パラメータ \( N \) および各偶数間隙 \( 2k \) に対して、ポリニャック 検査はいかなる尾部も \( g_n = 2k \) を回避しないことを検証する。任意の固定コンパクト化パラメータ \( N \) に対してアルゴリズムは停止する。連続極限 \( N \to \infty \) は古典的主張を回復する。

これで離散二重時空の枠組みにおける ゴールドバッハ 予想、ポリニャック 予想、および コラッツ 予想の証明を完了する。

\section{統一ディオファントス枠組みと降下}

先行各章の証明は共通の代数構造を共有する。我々はこの構造を一般枠組みへ抽象化する。以下の提示は、ダガーそれ自体が \(J\) を厳密に減少させるという以前の非形式的主張に対する、Lean~4 で機械検証された訂正を取り込む。

\subsection{ディオファントス方程式の一般表示}

第5--9章のディオファントス論証は\emph{冪和}方程式
\[
\sum_i c_i\, b_i^{e_i}=0
\]
（整数係数 \(c_i\)、底 \(b_i\)、指数 \(e_i\)）によって運ばれ、これは フェルマー、ビール、abc、および ゴールドバッハ 型主張の加法骨格を覆う。第4章の埋め込みを用いて、加法内容は零トランスレータ
\[
T\Bigl(\sum_i c_i\, b_i^{e_i}\Bigr)
\]
によって実現され、したがって方程式が消えるのはこのトランスレータが恒等元に等しいとき、かつそのときに限る（\(T\) の忠実性）。非零の底の間の乗法不整合はローター \(R(b_i)\) によって別に記録され、冪写像のもとで増幅されうる。非形式的な模式恒等式
\[
R(f) \cdot \exp\Bigl( \sum_i c_i T(x_i) \Bigr) = 1
\]
はしたがって、任意の多項式 AST の単一モーターへの字義どおりの埋め込みとしてではなく、この二セクター符号化として理解される。離散側では、次いで有限な許容集合を零高さ構成について探索するか、定理~\ref{thm:torsion-bound} に対する増幅によって矛盾を導く。（連続な整数ローターは量子化なしには許容格子点と\emph{同一視されない}。大きな \(|n|\) はすでに \(|\Jnorm(R(n))|>1\) を与える。）

\subsection{双対抽出対有限探索}
\label{sec:unified-descent}

\begin{proposition}[ダガーはねじれ高さを保存する]
\label{prop:dagger-height}
\(\Omega\) をパラメータ \((\alpha_a,\beta_a)\) をもつねじれ不整合とする。パラメータ \((\beta_a,\alpha_a)\) をもつ構成を \(\Omega^\dagger\) と書くと（パラメータ水準のダガー / 双対再構成）、
\[
J(\Omega^\dagger)=-J(\Omega),\qquad |\Jnorm(\Omega^\dagger)|=|\Jnorm(\Omega)|
\]
である。とくに、ダガーは正規化高さの厳密な減少を決して与えない。
\end{proposition}

\begin{proof}
定義~\ref{def:Jnorm} と
\[
J=\frac12\sum_a(\alpha_a^2-\beta_a^2)
\]
から直ちに従う。この恒等式は Lean~4 モジュール \texttt{DstDiophantine.\allowbreak Framework.\allowbreak Descent} において機械検証済みである。
\end{proof}

したがって、より古い物語「ダガー \(\to\) 単数群フィルタ \(\to\) 双対再構成が \(J\) を厳密に減少させる」は、述べられたとおりには正しくない。離散トーラスの単数群像は許容ローター代表元を記録するが、その像の現在の定義のもとでは、独立した高さ減少フィルタを課さない。

妥当なままであるもの——そして証明が実際に用いるもの——は次の三層スキームである。

\begin{enumerate}[label=(\roman*)]
  \item \textbf{忠実な整数符号化。} 冪和の消滅は対応する零トランスレータが \(1\) に等しいことと同値である。乗法不整合は増幅のために利用可能である。
  \item \textbf{有限な許容探索。} 各固定 \(N\) 上で、許容離散構成の集合は有限である。零の正規化高さはラピディティ添字上の整数格子条件と同値であり、したがって決定可能である。網羅的探索は有限ステップの後に停止する（同値に：有限な残存集合から候補を除去する任意のスキーマは整礎である）。
  \item \textbf{増幅 / 矛盾。} 特定の予想に対して、仮想の非自明解は、冪写像の後に \(|\Jnorm|>1\) となる構成を生じ、定理~\ref{thm:torsion-bound} に矛盾する。
\end{enumerate}

したがって停止は、許容トーラスの有限性（層~(ii)）または上限それ自体（層~(iii)）によって保証され、ダガー駆動の高さ降下によっては保証されない。ダガーは、高さを保存しつつ通常セクターと双対セクターを交換する双対対合として有用なままである。

\subsection{古典一般相対性理論の回復}

連続極限 \( N \to \infty \) において離散トーラスはコンパクト群 \( \mathrm{Spin}^+(3,1) \oplus \mathrm{Spin}^+(3,1) \) において稠密となる。有限離散解析は、許容セクターにおいて、古典変分原理
\[
\delta \int J \, d^4x = 0
\]
へ収束する。これはまさに一般相対性理論の遠平行等価理論（TEGR）である。したがってディオファントス論証を組織する同一の代数的対象が、いかなる追加構造なしに古典重力をも再現する。離散理論はしたがって数論的基礎と一般相対性理論の紫外完備化の双方を供給する。

この統一枠組みは、離散性と射影幾何構造を備えた双四元数代数が、算術と重力物理学の双方に対する単一の整合した設定として働くことを示す。

\section{含意と予測}

離散双四元数枠組みは、古典一般相対性理論および他の量子重力アプローチの双方からそれを区別する一連の物理的および数論的含意をもたらす。

\subsection{自然な紫外カットオフと特異点の解消}

厳密な上限 \( |\Jnorm| \leq 1 \) と離散二重時空の有限トーラスは、プランク尺度における自然な紫外カットオフを課す。曲率特異点は、反発層（\( J < 0 \)）が真の特異点が形成される前に崩壊を止める、有限かつ多層のねじれ構成によって置き換えられる。したがってブラックホール特異点と宇宙論的特異点の双方が、異質なプランク尺度物理学を呼び出すことなく解消される。

\subsection{離散ねじれスペクトル}

重力波および核励起は離散的なねじれ寄与を獲得する。連星合体の後、残骸は引力層と反発層が交替する成層内部をもつ。重力波は各ねじれ界面で部分反射し、時間遅延が
\[
\Delta t \sim \frac{GM}{c^3} \approx 10^{-5} \Bigl( \frac{M}{10 M_\odot} \Bigr) \text{ s}
\]
である特徴的なエコー列を生じる。核エネルギー準位も同様に、コンパクト化尺度 \(\ell_N = 2\pi / N\) に対して位数 \( \ell_N^{-1} \) の離散補正を受ける。

\subsection{安定核の有限個数}

各核子は三層ねじれ星として埋め込まれる。標準陽子および中性子ローター類 \( [R_{\rm p}], [R_{\rm n}] \in \mathrm{Spin}^+(3,1) \oplus \mathrm{Spin}^+(3,1) \) は、粒子安定性に関する 伴う論文で展開された層数え分類に従い、バリオン数 \(B=1\) および双対位相巻き数 \(w=3\) を運ぶ。質量数 \(A\) の原子核は順序積
\[
R_N(A) := \prod_{j=1}^{A} R_{\nu_j}
\]
によって表される。ここで各因子 \(R_{\nu_j}\) は \([R_{\rm p}]\) または \([R_{\rm n}]\) のいずれかの代表元であり、加法的質量数は零トランスレータ
\[
T(A) := \sum_{j=1}^{A} T(1)
\]
によって記録される。ここで \(T(1)\) は第4章の単位零増分である。

\begin{definition}
核構成 \(R_N(A)\) が\emph{代数的に安定}であるとは、
\begin{enumerate}[label=(\roman*)]
  \item 積 \(R_N(A)\) が離散双四元数環の単数群に属する（第4章の有限同値関係を除いて）こと、
  \item 双対位相が閉じて閉じたねじれループをなすこと、すなわち巻き数が、三層核子トポロジーと両立する整数層数 \(\ell\) に対して \(w = 3\ell\) を満たすこと、および
  \item 付随する不整合が \(|\Jnorm(R_N(A))| \leq 1\) を満たすこと
\end{enumerate}
である。
\end{definition}

\begin{theorem}[有限核図表]
各固定コンパクト化パラメータ \(N\) に対して、有限個の質量数 \(A\) のみが代数的に安定な核構成を許す。したがって安定核の質量数に対する有限上界
\[
A_{\max}(N) < \infty
\]
が存在する。
\end{theorem}

\begin{proof}
有限単数群に関する補題（付録）により、離散トーラス \((\mathbb{Z}/N\mathbb{Z})^6\) は有限個の相異なるローター類のみを与える。各代数的に安定な構成はこの有限集合に属さなければならない。核子類 \([R_{\rm p}]\) および \([R_{\rm n}]\) は固定されているから、写像 \(A \mapsto R_N(A)\) は、繰り返しまたは閉ループ条件 (ii) の違反の前に、高々有限個の相異なる単数群元しか生み出せない。\(|\Jnorm| > 1\) となるいかなる構成も第3章のねじれ有界性定理によって排除される。したがって安定質量数の集合は有限であり、\(A_{\max}(N) < \infty\) である。
\end{proof}

定理は有限性を厳密に確立する。それはそれ自体ではカットオフの数値を固定しない。物理的終端尺度を見積もるため、双対セクターが周期 \(4\pi\) でコンパクトであること、および原子核が微視的ねじれ星であることに注意する。バリオン密度が増すにつれ、\(J\) の逐次的な符号反転が反発層と引力層を蓄積する。層状振幅が上限 \(|\Jnorm| = 1\) に達すると、最外の反発殻（\(J < 0\)）が核子を放出する自発的ねじれバウンスを引き起こす。この最初のバウンスが起こる臨界質量数を \(A_*\) と記す。核飽和密度における層予算スケーリングと、\(A \sim 10^2\) 範囲における魔法数系列の経験的終端とをあわせると、桁の見積もり
\[
A_* = O(10^2), \qquad A_* \sim 300
\]
が得られる。この見積もりは定理定数ではなく、離散代数から導出されてもいない。それは \(4\pi\)-コンパクト双対幾何と経験的核図表から推論された未導出の発見的主張である。上記の定理が与えるのは \(A_{\max}(N)<\infty\) のみである。\(A > A_*\) に対して、発見的バウンス描像は安定な閉じたねじれループを排除するであろうが、その数値的カットオフは Lean 検証済みの上限ではない。

これらの含意は、離散二重時空理論が単なる数学的再定式化ではなく、算術と重力の双方に対する具体的帰結をもつ物理的に予測的な枠組みであることを示す。

\section{結論}

我々は、古典重力と算術数論が単一の32次元構造——離散双四元数二重時空代数——から現れる完全な代数的枠組みを構成した。双対時空のラピディティを有限トーラス \( (\mathbb{Z}/N\mathbb{Z})^6 \) へ離散化し、代数を射影幾何代数 \( G(3,1,1) \) へ埋め込むことにより、次の主要な結果を達成した。

第一に、ねじれ不整合スカラー \( J \) は許容構成上で \(|\Jnorm| \leq 1\) を満たし（定理~\ref{thm:torsion-bound}）、\(|J| = 3\pi^2/8\) における自然な紫外カットオフを供給する。

第二に、すべての整数は標準ローターおよび零トランスレータを通じて代数へ忠実に埋め込まれた。有限な離散ローター像はすべてのディオファントス方程式を有界格子上の有限探索へ変換する。第5章はフェルマーの二軸座席を記録する。混合種は加法軸を時間的零生成子へ漏らし、\(N_1\) の有限サンドイッチは楕円的であり、時間的零剰余は種が混合であるときに限り非零である。混合モーター残差は未閉鎖であり、古典的なフェルマーの最終定理は主張しない。第6--10章は次に対する共通の増幅プログラムを記録する：
\begin{itemize}
  \item ビール 予想、
  \item abc 予想、
  \item リーマン 仮説、
  \item ゴールドバッハ 予想、
  \item ポリニャック 予想（場合 \( 2k = 2 \) としての双子素数予想を含む）、および
  \item コラッツ 予想。
\end{itemize}
共有される増幅機構は、代数関係を満たし同時に上限 \( |\Jnorm| \leq 1 \) を破る構成が存在しないことと、許容離散探索空間の有限性（第10章）である。問題固有の橋渡しは未証明のままである。それら予想の無条件の主張は述べない。ダガーによる双対抽出は \(|\Jnorm|\) を保存し、それ自体は高さを減らす降下ではない。

第三に、本枠組みは連続極限 \( N \to \infty \) において一般相対性理論の遠平行等価理論（TEGR）を正確に回復する一方、離散構造は数論的基礎と古典重力の紫外完備化の双方を供給する。したがって同一の代数的対象が、追加の場、次元、あるいは量子化の要請なしに算術と重力を統一する。

本理論は具体的な物理的含意をもたらす。離散的な重力波エコー、各 \(N\) における有限個の許容核ローター類、および核エネルギー準位への離散的なねじれ補正である。数値的カットオフ \(A_* \lesssim 300\) は未導出の発見的主張のままである。

今後の課題には、整数双四元数環の単数の完全分類（現在の離散トーラスの像を超えて）、大きなコンパクト化パラメータに対する系統的数値検証、連続な整数ローターと許容格子点の間のより緊密な量子化の橋渡し、および本古典枠組みとねじれ量子重力に関する 伴う論文で展開された量子セクターとの完全な統合が含まれる。したがって離散双四元数代数は、数論と重力物理学の双方に対する統一的基礎の候補として立つ。

この再定式化において、連続体は有用な古典近似であった。最も深い水準における実在は、離散的、粒子局所的、かつ二重時間的である。時空は物質を曲げない。物質は自身の内在的双対時空を不整合にし、この不整合——集団的に知覚されるとき——が我々が重力と呼んできたものである。

\appendix

\section{小さなコンパクト化パラメータ \texorpdfstring{\( N \)}{N} に対する \texorpdfstring{\( J \)}{J} 値の明示的格子列挙}

離散トーラス上で定理~\ref{thm:torsion-bound} を例示するため、小さな \( N \) に対する\emph{許容}不整合構成（定義~\ref{def:admissible}）を列挙する。各 \( N \) に対して、六つのラピディティパラメータ \( (\omega_1,\omega_2,\omega_3,\phi_1,\phi_2,\phi_3) \) は有限集合
\[
\Bigl\{ \frac{2\pi k}{N} \Bigm| k = 0,1,\dots,N-1 \Bigr\}^6
\]
を渡り、\(\alpha_a,\beta_a \geq 0\) かつ \(\alpha_a + \beta_a \leq \pi/2\) を満たす点のみを残す。ねじれスカラー
\[
J = \frac12 \sum_{a=1}^3 (\alpha_a^2 - \beta_a^2), \qquad
\Jnorm = \frac{8}{3\pi^2} J
\]
をこの許容部分集合上で評価する。非許容トーラス点は \(|J| > \Jbound\) を満たしうる。

\subsection{場合 \texorpdfstring{$N = 4$}{N = 4}}

ここで \(4\mid N\) かつ \(\lfloor N/4\rfloor=1\) であるから、鋭い格子上限は \(|\Jnorm|\le 1\) である。二つの極値格子点 \(n_a=1\)、\(m_a=0\) および \(n_a=0\)、\(m_a=1\)（すべての軸）は許容であり、\(\Jnorm=\pm 1\)、同値に \(J=\pm\Jbound\) を達成する。これは一般格子定理の \(N=4\) の場合であり、すべての \(4^6=4096\) 個のトーラス点の独立な列挙ではない。

\subsection{場合 \texorpdfstring{$N = 8$}{N = 8}}

再び \(4\mid N\) であるから、\(\lfloor N/4\rfloor=2\) であり、鋭い上限は \(|\Jnorm|\le 1\) へ崩壊し、\(n_a=2\)、\(m_a=0\) および \(n_a=0\)、\(m_a=2\) で達成される。\(8^6=262144\) 個のトーラス点は列挙しない。格子評価はすでにすべての許容構成を覆う。

\subsection{一般的パターン}

任意の \( N \) に対して、許容格子上限は \(|\Jnorm|\le(4\lfloor N/4\rfloor/N)^2\) であり、等号 \(\pm 1\) が成立するのは \(4\mid N\) のとき、かつ二つの一様極値点においてのみである。\(4\nmid N\) ならば不等式は厳密である。\( N \) が増すにつれ、許容点上で到達可能な \(\Jnorm\) 値の離散集合は \([-1,1]\) において稠密となる。ディリクレ 核の層共鳴はコンパクト格子上の同期を制御する（伴う論文）。それらは初等的な許容上限を補うが、置き換えない。

これらの格子恒等式は Lean~4（\texttt{DstDiophantine.Framework.Lattice}）において機械検証済みであり、定理~\ref{thm:torsion-bound} の離散検査として力任せの列挙を置き換える。

\section[ベイカー–キャンベル–ハウスドルフ 展開と Spin 被覆]{ベイカー–キャンベル–ハウスドルフ 展開\\ と Spin 被覆の詳細}

\subsection{ローター冪に対する ベイカー–キャンベル–ハウスドルフ 公式}

ビール 予想、abc 予想、およびねじれ不整合の増幅プログラムにおいて、我々は繰り返しローターを整数冪へ上げる。\(\Omega = \exp(X)\) とし、ここで \( X \in \mathfrak{so}(3,1) \) は不整合二重ベクトルである。このとき
\[
\Omega^p = \exp(pX + \tfrac{p(p-1)}{2}[X,X] + \tfrac{p(p-1)(2p-1)}{12}[X,[X,X]] + \cdots).
\]
二次で打ち切ると（小さな不整合、または高次交換子が離散格子によって抑制されるときに妥当）、
\[
\log(\Omega^p) = pX + \tfrac{p(p-1)}{2}[X,X] + O(p^3)
\]
を得る。\( J \) を定義する二次形式へ代入すると、主要な増幅
\[
J(\Omega^p) = p^2 J(\Omega) + O(p^3 \cdot \text{higher nested commutators})
\]
が生じる。\( p \geq 3 \) に対して係数 \( p^2 \) は有限トーラス上の任意の高次補正を厳密に支配する。これは第5章に記録した単軸診断の増幅であり、二軸フェルマー座席を覆わない。ローターを混合モーター \( M=RT \) に置き換えても事情は変わらない。トランスレータはねじれローターによって正規化されるアーベル正規部分群をなすから \( M^p = R^p\,T_p \)（\( T_p \) はトランスレータ）であり、\( J \) の増幅はもっぱら \( R^p \) が担い、並進セクターがねじれスカラーへ逆流することはない。

\subsection{Spin 被覆と主枝制限}

群 \( \mathrm{Spin}^+(3,1) \) は固有順時 ローレンツ群の二重被覆である。したがって指数写像
\[
\exp : \mathfrak{so}(3,1) \to \mathrm{Spin}^+(3,1)
\]
は局所微分同相であるが大域的には一対一ではない。核は元 \( \pm 1 \) からなる。したがって対数の主枝に注意を制限する。それは開集合
\[
U = \{ g \in \mathrm{Spin}^+(3,1) \mid \| \log g \| < 2\pi \}
\]
上で定義され、ノルムはキリング形式によって誘導される。この近傍の内部で対数は一価かつ解析的である。

双対写像 \( X \mapsto Xi \) は不整合二重ベクトルの係数が主枝内部で
\[
\alpha_a + \beta_a = \delta_a, \qquad |\delta_a| \leq \frac{\pi}{2}
\]
を満たすことを強制する。この上限は、双対セクターにおける完全な \( 2\pi \) 回転が Spin 群の元 \( -1 \) に対応し、それが主領域の境界上にあるという事実から従う。より大きな偏差は分岐切断を横切ることを要求し、隣接粒子が連続な零トランスレータによってのみ異なるという大域的整合条件によって排除される。

\subsection{離散トーラス上の打ち切りの正当化}

有限トーラス \( (\mathbb{Z}/N\mathbb{Z})^6 \) 上ですべてのラピディティは \( 2\pi \) で押さえられる。したがって不整合二重ベクトル \( X \) は \( \|X\| \leq 3\pi \) を満たす。任意の固定コンパクト化パラメータ \( N \) に対して、ベイカー–キャンベル–ハウスドルフ 級数の高次項は一様に有界である。二次打ち切りによって生じる誤差はしたがって \( N \) のみに依存する定数によって制御され、連続極限で消える。これはディオファントス証明で用いられるすべての近似を正当化する。

ベイカー–キャンベル–ハウスドルフ 展開と Spin 被覆制限の組み合わせは、第6--7章で用いられる増幅論証を支える。初等的な許容上限 \(|\Jnorm| \leq 1\)（定理~\ref{thm:torsion-bound}）は BCH 打ち切りから独立である。

\section{補助補題の証明}

\subsection{離散ローター像に関する補題}

\begin{lemma}
ねじれローター写像のもとでの離散トーラス \( (\mathbb{Z}/N\mathbb{Z})^6 \) の像は有限である。
\end{lemma}

\begin{proof}
離散トーラス \( (\mathbb{Z}/N\mathbb{Z})^6 \) は有限集合である。有限集合の像は有限である。この対象は離散ローター像であり、整数双四元数次数の単数群ではない。その次数はここでは構成されず、ローレンツ型整数次数は無限個の単数をもちうる。
\end{proof}

\subsection{ディリクレ 核の除去可能特異点に関する補題}

\begin{lemma}
任意の整数 \( N \geq 2 \) および任意の整数 \( k \) に対して、ディリクレ 核は
\[
\lim_{\delta\theta \to 2\pi k} D_N(\delta\theta) = \pm 1
\]
を満たす。
\end{lemma}

\begin{proof}
ディリクレ 核は
\[
D_N(\theta) = \frac{\sin(N\theta/2)}{N\sin(\theta/2)}
\]
によって与えられる。\(\theta = 2\pi k\) において分子と分母の双方が消え、\( 0/0 \) の不定形を与える。L'Hôpital の法則を一度適用すると
\[
\lim_{\theta \to 2\pi k} D_N(\theta) = \frac{(N/2)\cos(N\theta/2)}{(N/2)\cos(\theta/2)} \Big|_{\theta=2\pi k} = \cos(N\pi k) = (\pm 1)^N = \pm 1
\]
を得、これは \( N \) に依存しない。したがってすべての見かけの特異点は除去可能である。
\end{proof}

\subsection{ローター冪のもとでのねじれ増幅に関する補題}

\begin{lemma}
\(\Omega = \exp(X)\) とし、\( X \in \mathfrak{so}(3,1) \) とする。このとき
\[
J(\Omega^p) = p^2 J(\Omega) + O(p^3)
\]
であり、含意される定数はコンパクト化パラメータ \( N \) のみに依存する。
\end{lemma}

\begin{proof}
ベイカー–キャンベル–ハウスドルフ 公式を二次まで適用する：
\[
\log(\Omega^p) = pX + \tfrac{p(p-1)}{2}[X,X] + O(p^3).
\]
\( J \) を定義する二次形式は次数2で斉次である。したがって主要項はちょうど \( p^2 J(X) \) を寄与する。すべての高次交換子は有限トーラス上で \( N \) のみに依存する定数によって有界である。ゆえに誤差項は \( O(p^3) \) である。
\end{proof}

\subsection{主枝上限に関する補題}

\begin{lemma}
\(\mathrm{Spin}^+(3,1) \oplus \mathrm{Spin}^+(3,1)\) 上の対数の主枝内部で、不整合係数は \( |\delta_a| \leq \pi/2 \) を満たす。
\end{lemma}

\begin{proof}
双対写像 \( X \mapsto Xi \) は、双対セクターにおける完全な \( 2\pi \) 回転を Spin 群の中心元 \( -1 \) と同一視する。主枝は対数のノルムが厳密に \( 2\pi \) より小さい開集合として定義される。いかなる偏差 \( |\delta_a| > \pi/2 \) も全不整合をこの近傍から退出させ、主枝の選択に矛盾する。したがって \( |\delta_a| \leq \pi/2 \) である。
\end{proof}

\subsection{根基--双対コンパクト性に関する補題}
\label{lem:radical-dual}

\begin{lemma}
\( (a,b,c) \) を \( a + b = c \) を満たす原始三つ組とする。\(\mathrm{rad}(abc)\) によって定まる双対セクター・ローター構成は、基数が \(\omega(\mathrm{rad}(abc))\) およびコンパクト化パラメータ \( N \) のみに依存する \(\mathrm{Spin}^+(3,1)\) の有限部分集合に属する。とくに、双対ラピディティは abc 不整合ローターの対数における各巡回係数 \( \phi_a \) に対して
\[
0 \leq \phi_a < \frac{2\pi}{N} \cdot \omega(\mathrm{rad}(abc))
\]
を満たす。
\end{lemma}

\begin{proof}
各素数 \( p \mid abc \) は双対ローター \( R_{\rm dual}(p) = \exp(2\pi \Gamma_p / N) \) を寄与し、その位相は離散トーラス上で \( 2\pi/N \) の整数倍である。\(\gcd(a,b,c) = 1\) であるから、\( a \)、\( b \)、および \( c \) の素因子は互いに素であり、全双対構成は \(\omega(\mathrm{rad}(abc))\) 個のそのような位相の積である。離散ローター像（離散ローター像に関する補題）はさらに積を許容類の有限集合へ制限する。したがって各巡回係数は、相異なる素因子の個数と格子間隔 \( 2\pi/N \) の積によって押さえられる。
\end{proof}

\subsection{品質--ねじれ不等式に関する補題}
\label{lem:quality-torsion}

\begin{lemma}
固定コンパクト化パラメータ \( N \) に対して定数 \( Q(N) > 1 \) が存在し、\( a + b = c \) かつ \( q(a,b,c) > Q(N) \) を満たすすべての原始三つ組 \( (a,b,c) \) は \( J(\Omega_{\rm abc}) > \Jbound \) を満たす。ここで \( \Omega_{\rm abc} \) は第7章の abc 不整合ローターである。
\end{lemma}

\begin{proof}
命題~\ref{prop:quality-torsion} により、大きな品質は abc ねじれ高さ \( H_{\rm abc} = J(\Omega_{\rm abc}) \) が \( c_1(N)(q - 1)^2 \) を超えることを強制する。\( Q(N) = 1 + (c_1(N) \varepsilon_N)^{-1/2} \) と選ぶ。ここで \(\varepsilon_N\) は離散トーラス上の最小非零ねじれ高さである。このとき \( q(a,b,c) > Q(N) \) は \( H_{\rm abc} > c_1(N) \varepsilon_N \geq \Jbound \) を含意する。ここで主枝上限および補題~\ref{lem:radical-dual} からの双対セクターのコンパクト性を用いる。
\end{proof}

\subsection{累積間隙蓄積に関する補題}
\label{lem:gap-accumulation}

\begin{lemma}
\(\{g_n\}\) を連続する素数間隙の列とし、偶数間隙 \( 2k \) を固定する。コンパクト化パラメータ \( N \geq 2 \) の離散トーラス上で
\[
S_M := \sum_{n=1}^{M-1} \frac{g_n - 2k}{p_n}
\]
を定義する。すべての \( n \geq N_0 \) に対して \( g_n \neq 2k \) ならば、すべての \( M \geq M_*(N,k,N_0) \) に対して \( S_M > 4\pi/N \) となる \( M_*(N,k,N_0) \) が存在し、したがって \(\Jnorm(\Sigma_M) > 1\) である。
\end{lemma}

\begin{proof}
\( n \geq N_0 \) かつ \( g_n \neq 2k \) に対して、各項 \( (g_n - 2k)/p_n \) は正の有理数であり、\( g_n \geq 2 \) かつ偶数であるため \( 2/p_n \) で下から押さえられる。部分和 \( S_M \) は \( M \to \infty \) として限りなく増大する。すべての \( M \geq M_*(N,k,N_0) \) に対して \( S_M > 4\pi/N \) となるよう \( M_*(N,k,N_0) \) を選ぶ。不整合二重ベクトル \(\Sigma_M\) は \( i\Gamma_a \) の張る空間に属し、キリング形式は主枝正規化を除いて \( J(\Sigma_M) = (1/16) S_M^2 \) を与える。十分大きな \( S_M \) に対してこれは \( J(\Sigma_M) > \Jbound \) を与え、したがって \(\Jnorm(\Sigma_M) > 1\) である。
\end{proof}

これらの補題は第3--9章のすべての論証に対する技術的基礎を供給する。

\end{document}
